\documentclass[11pt,oneside,openany]{book}
\usepackage[T1]{fontenc}
\usepackage{lmodern}
\usepackage[letterpaper,margin=1.10in,includeheadfoot]{geometry}
\usepackage{amsmath,amssymb,amsthm,mathtools}
\usepackage{microtype,xcolor,titlesec,fancyhdr,etoolbox,needspace}
\usepackage{hyperref}

\ifdefined\pdftrailerid
  \pdftrailerid{MathGenome Topics in Lie Theory}
\fi
\definecolor{ink}{HTML}{172923}
\definecolor{deepgreen}{HTML}{2B5242}
\definecolor{rulegreen}{HTML}{8A9F92}
\hypersetup{colorlinks=true,linkcolor=deepgreen,urlcolor=deepgreen,citecolor=deepgreen,
  pdftitle={Topics in Lie Theory},pdfauthor={MathGenome},
  pdfsubject={ALG-08-D01 through ALG-08-D03},bookmarksnumbered=true}
\setcounter{tocdepth}{1}
\setlength{\headheight}{15pt}
\setlength{\emergencystretch}{2em}
\newcommand{\runningchapter}{}
\raggedbottom
\titleformat{\chapter}[display]
  {\normalfont\color{ink}}{\color{deepgreen}\large\scshape Chapter \thechapter}{8pt}
  {\huge\bfseries}
\titlespacing*{\chapter}{0pt}{7pt}{22pt}
\titleformat{\section}{\normalfont\Large\bfseries\color{ink}}{\thesection}{0.7em}{}
\pagestyle{fancy}
\fancyhf{}
\fancyhead[L]{\small\sffamily\color{deepgreen}MathGenome}
\fancyhead[R]{\small\nouppercase{\runningchapter}}
\fancyfoot[C]{\small\thepage}
\renewcommand{\headrulewidth}{0.35pt}
\renewcommand{\chaptermark}[1]{\renewcommand{\runningchapter}{\chaptername\ \thechapter}}
\fancypagestyle{plain}{\fancyhf{}\fancyfoot[C]{\small\thepage}
  \renewcommand{\headrulewidth}{0pt}}

\newtheorem{theorem}{Theorem}[chapter]
\newtheorem{lemma}[theorem]{Lemma}
\newtheorem{proposition}[theorem]{Proposition}
\newtheorem{corollary}[theorem]{Corollary}
\theoremstyle{definition}
\newtheorem{definition}[theorem]{Definition}
\newtheorem{example}[theorem]{Example}
\theoremstyle{remark}
\newtheorem{remark}[theorem]{Remark}
\newcommand*\keeptheoremheading[1]{\BeforeBeginEnvironment{#1}{\Needspace{4\baselineskip}}}
\forcsvlist{\keeptheoremheading}{theorem,lemma,proposition,corollary,definition,example,remark}
\pretocmd{\section}{\Needspace{7\baselineskip}}{}{}
\allowdisplaybreaks[1]

\begin{document}
\frontmatter
\hypersetup{pageanchor=false}
\begin{titlepage}
\thispagestyle{empty}
\thispdfpagelabel{Cover}
\begin{flushleft}
\vspace*{0.72in}
{\large\scshape\color{deepgreen}MathGenome\par}
\vspace{0.36in}
{\color{rulegreen}\rule{0.94\textwidth}{0.65pt}\par}
\vspace{0.46in}
{\color{ink}\fontsize{35}{41}\selectfont\bfseries
  Topics in Lie Theory\par}
\vspace{0.34in}
{\color{deepgreen}\Large Infinite-dimensional, super, and restricted Lie algebras,\\[4pt]
  quantum groups, and deformations\par}
\vfill
{\normalsize ALG-08-D01 through ALG-08-D03\par}
\vspace*{0.50in}
\end{flushleft}
\end{titlepage}
\hypersetup{pageanchor=true}
\begingroup
\small
\makeatletter
\patchcmd{\l@chapter}{\vskip 1.0em \@plus \p@}
  {\vskip 0.45em \@plus 0.25\p@}{}{}
\makeatother
\tableofcontents
\endgroup

\chapter*{Scope and conventions}
\phantomsection
\addcontentsline{toc}{chapter}{Scope and conventions}
\renewcommand{\runningchapter}{Scope and conventions}

These notes continue \emph{Lie Algebras and Finite-Dimensional
Representations}. The first chapter constructs infinite-dimensional
Lie algebras from generalized Cartan data and studies the Witt and
Virasoro algebras. The affine loop model and the classification of its
integrable highest weights are supplied by the specified results of
\emph{Topics in Representation Theory}.

The second chapter treats the sign rules and enveloping algebra of a
Lie superalgebra, then the distinct restrictions that arise in
characteristic \(p\). It uses examples in which the defining relations
can be checked explicitly. The third chapter studies the quantized
enveloping algebra in type \(A_2\) and its relation to the classical
algebra. The fourth chapter computes obstruction classes for a Lie
bracket and for a deformation of the coproduct.

Each chapter states its base field and characteristic. We use the
coefficient conventions of the earlier MathGenome notes for each
numbered prerequisite and state sign conventions when passing to a
super bracket, a Lie bialgebra, or a skew coproduct. The source map
locates each numbered result and records the prerequisites of its
proof.

\mainmatter
\chapter{Contragredient Algebras and Virasoro}

All vector spaces in this chapter are over \(\mathbb C\).
The contragredient construction starts from a matrix and retains the
Cartan subalgebra while discarding the largest ideal disjoint from it.
This gives the structural meaning of the simple-root relations beyond
finite type. The second half constructs the Witt and Virasoro algebras
and computes their first highest-weight singular vectors.
For the loop bracket, the degree derivation and the concrete untwisted
affine highest-weight results we use \emph{Topics in Representation
Theory}, LO09--LO14 \cite[PDF pp.~26--28]{mg-reptopics}.

\section{Cartan data and the contragredient quotient}

\begin{definition}[Generalized Cartan matrices and realizations]
\label{res:KM01}
A generalized Cartan matrix \(A=(a_{ij})_{1\le i,j\le n}\) has
integer entries, \(a_{ii}=2\), \(a_{ij}\le0\) for \(i\ne j\), and
\(a_{ij}=0\) if and only if \(a_{ji}=0\). It is symmetrizable
if there are positive real numbers \(d_i\) such that
\(d_i a_{ij}=d_j a_{ji}\). A realization is a vector space
\(\mathfrak h\) with independent elements \(h_i\) and independent
linear forms \(\alpha_i\) such that
\(\alpha_j(h_i)=a_{ij}\). It is minimal when
\(\dim\mathfrak h=2n-\operatorname{rank}A\).
\end{definition}

\begin{lemma}[Existence and uniqueness of a minimal realization]
\label{res:KM02}
Every generalized Cartan matrix has a minimal realization. Any
two such realizations are related by an isomorphism carrying each
\(h_i\) to its counterpart and pulling back each counterpart of
\(\alpha_i\) to \(\alpha_i\). In a minimal realization,
\(\bigcap_i\ker\alpha_i\) has dimension \(n-\operatorname{rank}A\).
\end{lemma}

\begin{proof}
Let \(\mathfrak h_0\) have basis \(h_i\), and define
\(T_0:\mathfrak h_0\to\mathbb C^n\) by
\(T_0(h_i)=(a_{i1},\ldots,a_{in})\).
Choose a complement \(E\) of its image in \(\mathbb C^n\),
set \(\mathfrak h=\mathfrak h_0\oplus E\), and let
\(T(h,e)=T_0(h)+e\). The coordinate forms of \(T\) are independent
because \(T\) is onto. In any realization, the restriction of
\(T=(\alpha_i)_i\) to \(\mathfrak h_0\) has rank
\(\operatorname{rank}A\); surjectivity forces
\(\dim(\mathfrak h/\mathfrak h_0)\ge n-\operatorname{rank}A\).
For a minimal realization, choose a linear section of \(T\) over the
same complementary space \(E\). Its image is disjoint from
\(\mathfrak h_0\) and complements it by dimension. Thus the above
model identifies the realizations. The dimension of the common
kernel is rank-nullity for \(T\).
\end{proof}

\begin{proposition}[The universal algebra before the quotient]
\label{res:KM03}
Let \(\widetilde{\mathfrak g}\) be generated by the vector space
\(\mathfrak h\) and symbols \(e_i,f_i\), with relations
\[
 \begin{gathered}
 [h,h']=0,\qquad [h,e_i]=\alpha_i(h)e_i,\qquad
 [h,f_i]=-\alpha_i(h)f_i,\\
 [e_i,f_j]=\delta_{ij}h_i .
 \end{gathered}
\]
The displayed copy of \(\mathfrak h\) embeds,
and there is a vector space decomposition
\[
 \widetilde{\mathfrak g}
   =\widetilde{\mathfrak n}^-\oplus\mathfrak h
      \oplus\widetilde{\mathfrak n}^+ ,
\]
where the Lie subalgebras generated by the \(f_i\) and \(e_i\)
respectively are free on those generators. With
\(Q=\bigoplus_i\mathbb Z\alpha_i\), the grading has
\(\deg h=0\), \(\deg e_i=\alpha_i\), and
\(\deg f_i=-\alpha_i\). Each homogeneous space is the
corresponding weight space for the adjoint action of \(\mathfrak h\).
\end{proposition}

\begin{proof}
We use A05 and A07 of \emph{Lie Algebras and Finite-Dimensional
Representations} for PBW and the embedding of a free Lie algebra
in the corresponding tensor algebra \cite[PDF pp.~8--9]{mg-lie}.
Take the tensor space on symbols \(F_i\) and its vacuum \(v_\lambda\)
for any \(\lambda\in\mathfrak h^*\). On a word \(w\) let
\(\gamma(w)\) be the sum of the \(\alpha_j\) labeling its letters.
Let \(f_i\) prefix \(F_i\) and let \(h\) act diagonally on \(w\) by
\((\lambda-\gamma(w))(h)\). Define \(e_i\) to kill the vacuum and
recursively set
\[
 e_i(F_j w)=F_j(e_i w)
               +\delta_{ij}(\lambda-\gamma(w))(h_i)\,w.
\]
Every nonzero term of \(e_i w\) has \(\gamma\)-degree
\(\gamma(w)-\alpha_i\), by induction on the length. The weight
shift gives the required \(h,e_i\) bracket, and the recursion
gives exactly the \(e_i,f_j\) bracket. Thus these operators realize
the relations. If an \(h\) acts as zero for all \(\lambda\), its value
on each vacuum is zero, so \(h=0\). The action of words in the
\(f_i\) on the vacuum is free associative prefixing. Restricting
that action to the free Lie subalgebra proves freeness of
\(\widetilde{\mathfrak n}^-\). Interchanging \(e_i,f_i\) and sending
\(h\) to \(-h\) preserves the relations, hence gives freeness of
\(\widetilde{\mathfrak n}^+\) as well.

By induction on bracket length the relation for \([e_i,f_j]\)
gives \([e_i,\widetilde{\mathfrak n}^-]\subseteq
\mathfrak h+\widetilde{\mathfrak n}^-\): expand the bracket with
a Lie word into two terms by Jacobi. The induction terms
bracket elements of \(\mathfrak h\) or \(\widetilde{\mathfrak n}^-\)
with smaller negative words. The interchange gives the opposite
containment for \([f_i,\widetilde{\mathfrak n}^+]\).
The sum of the three subspaces contains the generators and is stable
under their adjoints. Iterated adjoints of the generators applied to
generators span the presented Lie algebra by Jacobi, so the sum
is all of it. The relations are homogeneous for \(Q\); the three
subspaces have only strictly negative, zero and strictly positive
degrees respectively. This gives directness. On a bracket word
the adjoint \(h\)-weight is its degree, by the derivation rule.
Independence of the \(\alpha_i\) identifies distinct degrees
with distinct weights.
\end{proof}

\begin{theorem}[The Cartan-disjoint ideal and the center]
\label{res:KM04}
There is a largest ideal \(\mathfrak r\) of
\(\widetilde{\mathfrak g}\) with \(\mathfrak r\cap\mathfrak h=0\).
The contragredient Lie algebra
\(\mathfrak g(A)=\widetilde{\mathfrak g}/\mathfrak r\) inherits
the triangular decomposition and the \(Q\)-grading.
Its center is \(\bigcap_i\ker\alpha_i\subset\mathfrak h\).
For a minimal realization its dimension is
\(n-\operatorname{rank}A\). Any ideal of \(\mathfrak g(A)\)
disjoint from \(\mathfrak h\) is zero.
\end{theorem}

\begin{proof}
Each ideal is a direct sum of its intersections with the weight
spaces. Indeed an element has only finitely many weight components.
Choose \(h\) separating their values, which is possible by avoiding
finitely many proper hyperplanes, and extract those components
as polynomial projections in \(\operatorname{ad}h\). Stability
under the adjoint keeps them in the ideal. Thus an ideal disjoint
from \(\mathfrak h\) has zero degree-zero component by
Proposition~\ref{res:KM03}. The sum of all these ideals still has
zero degree-zero component, and is the required largest ideal.
In the quotient, the triangular summands survive as the corresponding
degree ranges, including all of \(\mathfrak h\). A central vector is
of degree zero because it commutes with \(\mathfrak h\).
The \(e_i,f_i\) in the quotient are nonzero since
\([e_i,f_i]=h_i\ne0\); hence a vector of \(\mathfrak h\) commutes
with them precisely when all the \(\alpha_i\) vanish on it.
Lemma~\ref{res:KM02} gives the dimension. Finally the preimage of a
Cartan-disjoint ideal of the quotient is itself Cartan-disjoint and
contains \(\mathfrak r\), so it equals \(\mathfrak r\).
\end{proof}

\section{Simple reflections and highest weights}

\begin{theorem}[Serre relations in the quotient]
\label{res:KM05}
For \(i\ne j\), the simple generators of \(\mathfrak g(A)\) satisfy
\[
 (\operatorname{ad}e_i)^{1-a_{ij}} e_j=0,\qquad
 (\operatorname{ad}f_i)^{1-a_{ij}} f_j=0 .
\]
Each \(\operatorname{ad}e_i\) and \(\operatorname{ad}f_i\) is locally
nilpotent on \(\mathfrak g(A)\).
\end{theorem}

\begin{proof}
In \(\widetilde{\mathfrak g}\) write
\(v_s=(\operatorname{ad}e_i)^s e_j\) and \(m=1-a_{ij}\).
One has \([f_i,v_0]=0\) and \(v_s\) has \(h_i\)-weight
\(a_{ij}+2s\). Thus induction gives
\[
 [f_i,v_s]=-s(a_{ij}+s-1)v_{s-1}\qquad(s\ge1),
\]
by expanding \([f_i,[e_i,v_{s-1}]]\) using Jacobi and
\([f_i,e_i]=-h_i\). At \(s=m\) this is zero.
For \(k\notin\{i,j\}\) the adjoint of \(f_k\) commutes with
\(\operatorname{ad}e_i\) and kills \(e_j\), so it kills \(v_m\).
The \(f_j\)-bracket is \(-(\operatorname{ad}e_i)^m h_j\).
It vanishes when \(m\ge2\), since the first bracket is already a
multiple of \(e_i\). When \(m=1\), \(a_{ij}=0\) gives
\(a_{ji}=0\), so it still vanishes.

Take the span of \(v_m\) and all successive adjoints of the
\(e_k\) on it. It is stable under \(\mathfrak h\) by homogeneity,
and under the \(e_k\) by construction. It is stable under each
\(\operatorname{ad}f_k\), by commuting this operator past the
first \(\operatorname{ad}e_\ell\) and using
\([f_k,e_\ell]=-\delta_{k\ell}h_k\); induction uses that \(f_k\)
kills \(v_m\). This span is therefore an ideal, with strictly
positive degrees, and lies in the ideal of
Theorem~\ref{res:KM04}. The interchange \(e_i\leftrightarrow f_i\),
\(h\mapsto-h\), preserves \(\mathfrak r\) by its maximality, and
gives the negative relations. On the simple generators, local
nilpotence of an adjoint \(\operatorname{ad}e_i\) follows from
these relations, from \([e_i,f_j]=\delta_{ij}h_i\), and from
\([e_i,h]=-\alpha_i(h)e_i\). The derivation rule makes it
locally nilpotent on every bracket word, because a sufficiently
high iterate then hits some generator beyond its bound.
The negative case follows by interchange.
\end{proof}

\begin{proposition}[Reflection automorphisms and real roots]
\label{res:KM06}
For each \(i\) there is an automorphism
\[
 r_i=\exp(\operatorname{ad}e_i)
      \exp(-\operatorname{ad}f_i)\exp(\operatorname{ad}e_i)
\]
whose restriction to \(\mathfrak h\) sends
\(h\) to \(h-\alpha_i(h)h_i\).
It takes the root space of \(\beta\) isomorphically to that
of \(s_i\beta=\beta-\beta(h_i)\alpha_i\). All roots in the
orbits of the simple roots under the \(s_i\) have multiplicity one.
If \(A\) is symmetrizable, the form on the real span of the
root lattice with
\((\alpha_i,\alpha_j)=d_i a_{ij}\) is invariant under each
\(s_i\). In this case a root in the orbit of \(\alpha_i\)
has squared length \(2d_i\), even when the form is degenerate.
\end{proposition}

\begin{proof}
All exponentials are finite sums on each vector by
Theorem~\ref{res:KM05}. The iterated derivation rule for brackets
shows \(\exp(D)[x,y]=[\exp(D)x,\exp(D)y]\); the inverse is
\(\exp(-D)\), by the binomial identity. In the three factors above,
the first applied to \(h\) gives \(h-\alpha_i(h)e_i\),
the second sends \(h\) to \(h-\alpha_i(h)f_i\) and \(e_i\) to
\(e_i+h_i-f_i\). Applying the last factor gives \(h-\alpha_i(h)h_i\),
using \([h_i,e_i]=2e_i\). This map on \(\mathfrak h\) is an
involution, so the new weight of \(r_i x\), for \(x\) of
weight \(\beta\), is \(\beta\) composed with that map. The
\(\alpha_i\) root space is the line \(\mathbb C e_i\): a bracket
word of that degree uses just \(e_i\), and \(e_i\ne0\).
Its orbit thus has the indicated multiplicities. The formula
\(\beta(h_i)= (\alpha_i,\beta)/d_i\) expresses \(s_i\) as
reflection in \(\alpha_i\) for the symmetric form. Direct
expansion preserves the form since
\((\alpha_i,\alpha_i)=2d_i\); no inverse of the form is used.
\end{proof}

\begin{example}[An infinite rank-two root orbit]
\label{res:KM07}
For \(a,b>0\) integers with \(ab\ge4\), let
\(A=\left(\begin{smallmatrix}2&-a\\-b&2\end{smallmatrix}\right)\).
The algebra \(\mathfrak g(A)\) is infinite-dimensional.
\end{example}

\begin{proof}
In simple-root coordinates the matrix of \(s_1s_2\) is
\(\left(\begin{smallmatrix}ab-1&-a\\ b&-1\end{smallmatrix}\right)\),
of determinant \(1\) and trace \(ab-2\). If \(ab>4\), its
two eigenvalues are positive reciprocal real numbers, different
from \(1\). Thus no nonzero vector is fixed by any positive
power of this matrix. If \(ab=4\), the matrix is \(I+N\) with
\(N^2=0\) by its characteristic polynomial, and
\(N\alpha_1=2\alpha_1+b\alpha_2\ne0\).
Then \((I+N)^j\alpha_1=\alpha_1+jN\alpha_1\) are all distinct.
In either case Proposition~\ref{res:KM06} gives infinitely many
nonzero, distinct root spaces in the orbit of \(\alpha_1\).
\end{proof}

\begin{theorem}[Highest-weight modules and integrability]
\label{res:KM08}
For \(\lambda\in\mathfrak h^*\), the Verma module induced from the
one-dimensional \(\mathfrak h\oplus\mathfrak n^+\)-module of
weight \(\lambda\) has finite-dimensional weight spaces,
one-dimensional top, and a unique simple quotient \(L(\lambda)\).
On this simple quotient all \(e_i,f_i\) are locally nilpotent
if and only if \(\lambda(h_i)\in\mathbb Z_{\ge0}\) for all \(i\).
In that case the multiplicities of its weight spaces are invariant
under the reflections \(s_i\) of Proposition~\ref{res:KM06}.
\end{theorem}

\begin{proof}
PBW, Result A05 of \cite[PDF pp.~8--9]{mg-lie}, for an order
with the negative-root basis first gives the
Verma module as \(U(\mathfrak n^-)\) applied freely to the
highest vector. Every nonzero degree in the negative algebra
is a strictly negative combination of the simple degrees.
For a given weight difference \(\gamma=\sum_i m_i\alpha_i\)
with \(m_i\ge0\), its negative enveloping component is spanned
by words in the \(f_i\) with exactly \(m_i\) occurrences of
each \(i\), of which there are finitely many.
The degree of the vacuum is zero, and every proper submodule
misses the vacuum. Such submodules are sums of weight components,
by the polynomial extraction used in Theorem~\ref{res:KM04}.
Their sum is therefore still proper and is the unique maximal
submodule.

On a highest vector \(v\) the rank-one relation gives
\[
 e_i f_i^r v
   =r(\lambda(h_i)-r+1)f_i^{r-1}v
\]
by induction, as in the affine criterion LO12
\cite[PDF p.~27]{mg-reptopics}. If \(r>0\) is minimal with
\(f_i^r v=0\), its coefficient must vanish, proving necessity.
Conversely, for \(m=\lambda(h_i)\in\mathbb Z_{\ge0}\),
the vector \(f_i^{m+1}v\) in the Verma module is killed by \(e_i\)
by the formula, and by every other \(e_j\) since \([e_j,f_i]=0\).
It generates a proper submodule: PBW shows that the generated
weights are at most its own lower weight. It therefore vanishes
in the simple quotient. Each \(e_i\) already kills its highest
vector. For \(x=e_i\) or \(f_i\), Theorem~\ref{res:KM05} and the
derivation rule give local nilpotence of \(\operatorname{ad}x\)
on \(U(\mathfrak g(A))\). Now
\[
 x^N(uv)=\sum_{j=0}^N\binom Nj
       (\operatorname{ad}x)^j(u)x^{N-j}v
\]
proves local nilpotence on the entire module by taking \(N\)
large enough for both bounds. On an integrable module the
products \(\exp(e_i)\exp(-f_i)\exp(e_i)\) act invertibly.
The same binomial formula shows that conjugation by
\(\exp(x)\) on the action of an algebra element is
the action of \(\exp(\operatorname{ad}x)\). In particular the
product intertwines the action of \(\mathfrak h\) with that
of its image under \(r_i\). It therefore takes the weight
space of \(\mu\) isomorphically to that of \(s_i\mu\).
\end{proof}

\section{Comparison with the untwisted affine model}

\begin{proposition}[The contragredient affine algebra is the loop model]
\label{res:KM09}
Let \(\mathfrak a\) be finite-dimensional complex simple and take
the untwisted affine loop algebra with its degree derivation
as in LO09--LO10. The simple forms
\(\alpha_0=\delta-\theta,\alpha_1,\ldots,\alpha_\ell\) and
coroots \(h_0=K-h_\theta,h_1,\ldots,h_\ell\) give a
minimal realization of a symmetrizable generalized Cartan
matrix. The associated \(\mathfrak g(A)\) is the full loop
model, with the center and the derivation kept distinct.
\end{proposition}

\begin{proof}
Use the construction and conventions of LO09--LO10
\cite[PDF p.~26]{mg-reptopics}. Here \(\theta\) is the highest
finite root. The simple forms are independent since
\(\alpha_0(d)=1\) and all finite simple forms vanish on \(d\).
The coroots are independent by the coefficient of \(K\).
The finite simple Cartan block is invertible. The full matrix
has rank \(\ell\): the nonzero element \(K\) in the span of the
coroots is annihilated by all simple forms, and the finite block
gives the lower bound. Thus the realization of dimension
\(\ell+2\) is minimal. In the entries involving \(0\) and a
finite simple index, the two off-diagonal entries are
\(-\alpha_i(h_\theta)\) and \(-\theta(h_i)\), both nonpositive
integers by the root strings and dominance of \(\theta\), and
both zero exactly when \((\theta,\alpha_i)=0\), using D04 and
E09 of \cite[PDF pp.~24--25, 34--35]{mg-lie}.
The \(0,0\) entry is \(2\). Extending the finite root form
by making \(\delta\) orthogonal to everything gives
\(a_{ij}=2(\alpha_i,\alpha_j)/(\alpha_i,\alpha_i)\).
Thus the positive half-squared lengths symmetrize the full matrix.

The simple triples and the affine Cartan subalgebra satisfy the
relations of Proposition~\ref{res:KM03}, and generate the loop
model by LO10. The resulting surjection is injective on
\(\mathfrak h\). To see that its kernel is the largest
Cartan-disjoint ideal, any nonzero ideal of the loop model
contains a nonzero weight vector by polynomial weight
extraction. A nonzero real root vector
\(x_\alpha\otimes t^m\) brackets with a suitable opposite
root vector in degree \(-m\) to give a nonzero Cartan
element: the finite component is a nonzero multiple of
\(h_\alpha\). An imaginary degree vector
\(h\otimes t^m\) with \(m\ne0\) brackets with some
\(h'\otimes t^{-m}\) to give
\(m B(h,h')K\ne0\), since the invariant form is nondegenerate
on the finite Cartan. A nonzero element of weight zero
already lies in the affine Cartan. The largest ideal
\(\mathfrak r\) of Theorem~\ref{res:KM04} contains the kernel,
and its image has zero intersection with the affine Cartan:
lift a purported Cartan element along the isomorphism on
\(\mathfrak h\) and subtract that lift to land in the
kernel. Thus the image of \(\mathfrak r\) is zero, proving
the comparison.
\end{proof}

\section{Witt and its universal central extension}

\begin{proposition}[The Witt algebra is simple and perfect]
\label{res:KM10}
Every derivation of \(\mathbb C[t,t^{-1}]\) is a finite linear
combination of the elements \(L_m=-t^{m+1}\partial_t\),
\(m\in\mathbb Z\). These form the Witt algebra \(W\) with bracket
\([L_m,L_n]=(m-n)L_{m+n}\). The Lie algebra \(W\) is simple
and equals its derived algebra.
\end{proposition}

\begin{proof}
A derivation is determined by its value at \(t\), because the
Leibniz rule gives its values on \(t^n\) for positive powers
and \(D(t^{-1})=-t^{-2}D(t)\). The commutator of the
displayed derivations is computed at \(t\), giving the bracket.
The eigenvalues \(-m\) of \(\operatorname{ad}L_0\) are distinct.
Every element has finite support, so any nonzero ideal contains
some \(L_m\) by polynomial weight extraction. If \(m\ne0\)
then \([L_m,L_{-m}]=2mL_0\) lies in the ideal. With \(L_0\)
in the ideal, all \(L_n\), \(n\ne0\), are obtained by brackets.
The same two bracket formulas show that \(W=[W,W]\).
\end{proof}

\begin{theorem}[The Virasoro cocycle and the universal extension]
\label{res:KM11}
For an arbitrary complex vector space \(Z\) with trivial \(W\)-action,
every central \(Z\)-valued Lie \(2\)-cocycle of \(W\), modulo
coboundaries, has a unique representative of the form
\[
 c(L_m,L_n)=\delta_{m,-n}\frac{m^3-m}{12}\,z,\qquad z\in Z.
\]
In particular \(H^2(W,\mathbb C)\) has dimension one.
The Lie algebra \(\operatorname{Vir}=W\oplus\mathbb C C\)
with bracket
\[
 [L_m,L_n]=(m-n)L_{m+n}
             +\delta_{m,-n}\frac{m^3-m}{12}C,\qquad [C,L_m]=0 ,
\]
is perfect, has center \(\mathbb C C\), and is the universal
central extension of \(W\).
\end{theorem}

\begin{proof}
For trivial coefficients a cocycle is an alternating bilinear map
satisfying
\[
 c([x,y],u)+c([y,u],x)+c([u,x],y)=0 .
\]
Replacing it by \(c-\phi\circ[-,-]\) corresponds to changing
the linear section of a central extension by \(\phi\).
For \(j\ne0\) choose
\(\phi(L_j)=-c(L_0,L_j)/j\), so that after the change
\(c(L_0,L_j)=0\) for every \(j\). The cocycle equation on
\((L_0,L_m,L_n)\) now gives
\((m+n)c(L_m,L_n)=0\). Set \(b_m=c(L_m,L_{-m})\).
A further change with \(\phi\) supported at \(L_0\) and
\(\phi(L_0)=b_1/2\) makes \(b_1=0\) and retains the first
normalization. Then \(b_{-m}=-b_m\). On
\((L_m,L_n,L_{-m-n})\) the cocycle equation is
\[
 (m-n)b_{m+n}-(m+2n)b_m+(2m+n)b_n=0 .
\]
With \(n=1\) and \(m\ge2\) this gives
\((m-1)b_{m+1}=(m+2)b_m\). The unique solution for \(m\ge2\)
is \(b_m=(m^3-m)b_2/6\); the normalized values for
\(m=0,\pm1\) and antisymmetry give all integers. Conversely
those values satisfy the three-term relation for every \(m,n\),
as follows by expanding the cubic polynomial, and the cocycle
equation is zero off total index zero. This proves existence
and the stated form with \(z=2b_2\). A coboundary preserving
the first normalization has \(\phi(L_j)=0\) for \(j\ne0\),
and preserving \(b_1=0\) forces \(\phi(L_0)=0\).
Thus the normalized representative is unique.

The displayed bracket on \(\operatorname{Vir}\) is therefore
Lie. Its center is exactly \(\mathbb C C\) by simplicity and
noncommutativity of \(W\). In \(\operatorname{Vir}\), all
\(L_m\) and \(C\) are brackets or linear combinations thereof:
use first the brackets of \(L_0\) with \(L_n\), then
\([L_1,L_{-1}]=2L_0\) and
\([L_2,L_{-2}]=4L_0+C/2\). For any central extension
\(E\to W\), choose a linear section and form the cocycle with
values in its central kernel \(Z\).
After normalization, the chosen lifts of \(L_m\) and the vector
\(z\in Z\) define a Lie homomorphism
\(\operatorname{Vir}\to E\) over \(W\) by sending \(C\) to \(z\).
Two such homomorphisms differ by a central-valued linear map
and so agree on every bracket. Perfectness implies their equality.
This is the asserted universal property.
\end{proof}

\begin{proposition}[Witt derivations lift to the affine central extension]
\label{res:KM12}
In the loop model from Proposition~\ref{res:KM09}, before adjoining
the degree derivation, write \(x_i=x\otimes t^i\).
The assignments
\[
 D_m(x_i)=-i x_{m+i},\qquad D_m(K)=0
\]
give an action of \(W\) by derivations on the central extension
of the loop algebra. They therefore give a semidirect Lie algebra
with \(\operatorname{Vir}\) acting through \(W\); in this action
\(D_0\) is the negative of the loop degree derivation.
\end{proposition}

\begin{proof}
For the loop part, \(D_m\) applies \(L_m\) to the Laurent-polynomial
factor, so it is a derivation and the commutators give the Witt
action of Proposition~\ref{res:KM10}. For \(x_i,y_j\), the central
terms on the right side of the derivation identity sum to
\[
 \bigl(-i(m+i)-ji\bigr)\delta_{m+i+j,0}B(x,y)K=0 .
\]
This uses exactly the old loop cocycle of LO09. The central
term on the left side also vanishes since \(D_m(K)=0\).
The commutator equalities extend to \(K\), where all operators
vanish. Composing with the quotient
\(\operatorname{Vir}\to W\) gives an action of \(\operatorname{Vir}\);
the semidirect bracket satisfies Jacobi by the derivation and
commutator identities. The degree formula is immediate from
\(D_0(x_i)=-i x_i\).
\end{proof}

\section{Virasoro highest weights}

\begin{theorem}[Virasoro Verma modules]
\label{res:KM13}
For \(c,h\in\mathbb C\), there is a universal cyclic
highest-weight module \(M(c,h)\) with vector \(v\) satisfying
\(Cv=cv\), \(L_0v=hv\), \(L_jv=0\) for \(j>0\).
It has a basis
\[
 L_{-i_1}\cdots L_{-i_s}v,\qquad
       1\le i_1\le\cdots\le i_s,\quad s\ge0 .
\]
The weight of this vector for \(L_0\) is \(h+i_1+\cdots+i_s\).
Consequently the graded dimension with the top in degree zero is
\(\prod_{i\ge1}(1-q^i)^{-1}\) as a formal power series,
and \(M(c,h)\) has a unique simple quotient.
\end{theorem}

\begin{proof}
Induce the one-dimensional module of
\(\mathbb C C\oplus\bigoplus_{j\ge0}\mathbb C L_j\)
with the specified action. The bracket relations of
Theorem~\ref{res:KM11} show that this one-dimensional action
is well-defined. Induction gives universality. PBW, Result A05 of
\cite[PDF pp.~8--9]{mg-lie}, for a total order with the negative
generators first gives the indicated basis.
The bracket \([L_0,L_{-i}]=iL_{-i}\) gives the weight.
Choosing the number of occurrences of each \(i\) gives the product
of geometric series, and each degree has only finitely many such
choices. A proper submodule misses the top since the top generates,
and it splits into weight spaces by polynomial extraction.
The sum of all proper submodules thus still misses the top
and gives the unique simple quotient.
\end{proof}

\begin{example}[Singular vectors in degrees one and two]
\label{res:KM14}
In \(M(c,h)\), a nonzero vector of degree \(1\) killed by all
\(L_j\) with \(j>0\) exists exactly when \(h=0\). In degree \(2\)
such a vector exists exactly when
\[
 18h-(4h+2)(4h+c/2)=0 .
\]
In that case the space of these vectors is the line spanned by
\[
 3 L_{-1}^{\,2}v-(4h+2)L_{-2}v .
\]
Both kinds of vectors generate proper submodules. In particular,
at \((c,h)=(1,1/4)\) the first such degree is \(2\), where the
vector is proportional to \(L_{-1}^{\,2}v-L_{-2}v\).
\end{example}

\begin{proof}
In degree \(1\) the basis is \(L_{-1}v\) and \(L_1L_{-1}v=2hv\).
Higher \(L_j\) kill this vector by their weight shifts.
In degree \(2\), take the ordered basis
\((L_{-2}v,L_{-1}^2v)\). The maps \(L_1,L_2\) to the
degrees \(1,0\), respectively, have rows
\[
 (3,4h+2), \qquad (4h+c/2,6h).
\]
For instance \([L_1,L_{-2}]=3L_{-1}\);
commuting \(L_1\) past both factors of \(L_{-1}^2\)
gives \(2(L_0L_{-1}+L_{-1}L_0)v=(4h+2)L_{-1}v\).
Similarly \([L_2,L_{-2}]=4L_0+C/2\) and commuting
\(L_2\) past both factors of \(L_{-1}^2\) gives \(6hv\).
The maps \(L_j\) for \(j\ge3\) again kill this space by degree.
The determinant of the two rows gives the criterion, and since
the first entry of the first row is \(3\), its nullspace is the
displayed line when the determinant is zero.
PBW moves all positive generators to the right of a word
applied to a singular vector, where they kill it, and moves
the Cartan factors to act as scalars. Thus the submodule generated
in either positive degree cannot contain the original top.
Substituting \(c=1,h=1/4\) gives the example.
\end{proof}

The contragredient construction and root symmetries are compared with
Kac, Chapter 1 and Sections 3.1--3.7 \cite{kac}.
The untwisted affine convention agrees with Hernandez,
Sections 3--4 \cite{hernandez-affine}. For Virasoro and highest
weights, see Kac, Sections 7.3 and 9.14 \cite{kac}.
The converse assertion that the Serre ideal already equals the
Cartan-disjoint ideal in symmetrizable type, the Weyl--Kac character
formula, and the higher-degree classification of Virasoro singular
vectors are further topics.

\chapter{Lie Superalgebras and Restricted Lie Algebras}

The two theories in this chapter modify different parts of the ordinary
Lie algebra formalism. A superalgebra changes the interchange law for
homogeneous elements. In positive characteristic an ordinary Lie
algebra can carry a \(p\)-operation that is compatible with the
\(p\)-th power in an associative algebra. We prove the corresponding
basis theorems for both enveloping algebras and use them in small
rank examples.

\section{Parity and super enveloping algebras}

\begin{definition}[Lie superalgebras and their modules]
\label{res:SR01}
Let \(k\) be a field of characteristic different from two.
A Lie superalgebra is a graded vector space
\(\mathfrak a=\mathfrak a_{\overline 0}\oplus\mathfrak a_{\overline 1}\)
with an even bilinear bracket, meaning that degrees add. For
homogeneous \(x,y,z\), write
\(\epsilon_{xy}=(-1)^{|x||y|}\).
The identities are
\[
 [x,y]=-\epsilon_{xy}[y,x],
 \qquad [x,[y,z]]=[[x,y],z]+\epsilon_{xy}[y,[x,z]].
\]
In characteristic three we also require
\([x,[x,x]]=0\) for each odd \(x\). In other allowed characteristics
this follows from the first two identities. A graded module is a
graded vector space \(M\) with maps \(x\mapsto \rho(x)\) carrying
parity \(|x|\), satisfying
\(\rho([x,y])=\rho(x)\rho(y)-\epsilon_{xy}\rho(y)\rho(x)\).
Morphisms preserve the grading.
\end{definition}

Throughout the examples involving complex representations, the
ground field is \(\mathbb C\). The parity and basis arguments
use only the hypotheses of Definition~\ref{res:SR01}.
For the superalgebra conventions see Musson
~\cite[Section 1.1, PDF pp.~22--24]{musson}
and Serganova~\cite[Section 2]{serganova}.

\begin{proposition}[The supercommutator and the smallest matrix example]
\label{res:SR02}
In an associative graded algebra the bracket
\([x,y]=xy-\epsilon_{xy}yx\) makes its graded vector space a
Lie superalgebra. In \(\mathfrak{gl}(1|1)\), with even
generators \(z=E_{11}+E_{22}\), \(h=E_{11}\) and odd
generators \(e=E_{12}\), \(f=E_{21}\), the nonzero brackets
up to graded skew symmetry are
\[
 [h,e]=e,\qquad [h,f]=-f,\qquad [e,f]=z .
\]
In particular \(z\) is central and \(e,f\) are odd with
zero self-bracket. The supertrace is a linear functional
on \(\mathfrak{gl}(1|1)\) vanishing on every bracket, with
kernel \(\operatorname{span}_k(z,e,f)\).
\end{proposition}

\begin{proof}
Skew symmetry follows by interchanging the two products.
Expanding the derivation identity for \([x,-]\) on a
product gives
\[
 [x,yz]=[x,y]z+\epsilon_{xy}y[x,z]:
\]
the intermediate \(yxz\) terms cancel. Apply this to both
products in \([y,z]\) and use
\(\epsilon_{x,yz}=\epsilon_{xy}\epsilon_{xz}\) to obtain
the graded Jacobi identity. If \(x\) is odd then
\([x,x]=2x^2\), which commutes with \(x\). This also
verifies the extra identity in characteristic three.
For an abstract odd \(x\), substituting \(x=y=z\) into
Jacobi and using skew symmetry gives \(3[x,[x,x]]=0\).

For matrices take block-diagonal elements to be even
and off-diagonal elements odd. Multiplication of the
matrix units gives the displayed brackets
and \(e^2=f^2=0\).
The supertrace sends \(\operatorname{diag}(a,b)\)
to \(a-b\) and sends odd elements to zero. The
bracket table shows directly that its kernel is
the asserted span and contains the bracket subspace.
\end{proof}

\begin{theorem}[Super Poincare--Birkhoff--Witt]
\label{res:SR03}
Let \(\mathfrak a\) satisfy Definition~\ref{res:SR01}, and
order a homogeneous basis. Its enveloping superalgebra is
the quotient of \(T_k(\mathfrak a)\) by
\[
 xy-\epsilon_{xy}yx-[x,y]\qquad(x,y\text{ homogeneous}).
\]
The weakly increasing words in this basis, with each
odd basis element appearing at most once, form a basis.
For the word-length filtration its associated graded is
\[
 \operatorname{gr} U(\mathfrak a)
   \cong S(\mathfrak a_{\overline 0})
              \otimes_k\bigwedge(\mathfrak a_{\overline 1})
\]
as graded superalgebras. Graded modules over the enveloping
algebra and graded Lie superalgebra modules agree.
\end{theorem}

\begin{proof}
For a descending pair \(a>b\) replace \(ab\) by
\(\epsilon_{ab}ba+[a,b]\). For an odd repeated letter
replace \(aa\) by \([a,a]/2\). Order reductions first
by length and then by number of inversions. Each
replacement has smaller terms for this order, so
recursive reduction terminates. The relations of the
presentation are generated by these replacements:
for a weakly increasing pair use graded skew symmetry
or the repeated-letter relation, and extend bilinearly.

Disjoint adjacent replacements commute. Two overlapping
replacements have input \(abc\) with \(a\ge b\ge c\),
where a repeated adjacent pair must be odd. If all
three are distinct, reorder the three-letter term to
\(\epsilon_{ab}\epsilon_{ac}\epsilon_{bc}cba\) in
both paths. The difference of the shorter terms reduces
to
\[
 [[a,b],c]+\epsilon_{ab}[b,[a,c]]-[a,[b,c]]=0.
\]
Indeed the respective pairs of shorter terms are
\([a,b]c-\epsilon_{ac}\epsilon_{bc}c[a,b]\),
\(\epsilon_{ab}b[a,c]-\epsilon_{bc}[a,c]b\), and
\(\epsilon_{ab}\epsilon_{ac}[b,c]a-a[b,c]\).
For input \(aac\) with \(a\) odd, the difference reduces to
\(\tfrac12[[a,a],c]-[a,[a,c]]=0\).
For \(abb\) with \(b\) odd, it reduces to
\([[a,b],b]-\tfrac12[a,[b,b]]=0\).
These equalities are Jacobi with repeated odd entries.
Finally for \(aaa\) with \(a\) odd the difference reduces to
\(\tfrac12[[a,a],a]=0\), using the odd triple identity.
The checks are valid for bracket values that are sums of
basis elements by linearity.

Induct on the length-inversion order to compare any two possible
first steps. In the disjoint case the two paths have the same
common expansion. In the overlap case the expansion comparison
above leaves terms of smaller length, whose reductions agree by
induction. Prefixes and suffixes preserve this smaller-length
comparison. This gives a unique linear normal form. It kills
each replacement relation with every prefix and suffix by
the just-proved independence of first step, and so kills
their two-sided ideal. Conversely each reduction differs
from its input by an element of that ideal. Thus the normal
words are a basis of the quotient. Their leading symbols
give the displayed symmetric and exterior factors:
the leading relations say that even letters commute and
odd letters anticommute and square to zero. The tensor
algebra action on a graded module descends exactly when it
satisfies the signed bracket identity in
Definition~\ref{res:SR01}, giving the module claim.
\end{proof}

\begin{example}[Simple graded modules for the complex \(\mathfrak{gl}(1|1)\)]
\label{res:SR04}
On a finite-dimensional simple graded module the central \(z\)
acts as a scalar \(\gamma\). If \(\gamma=0\), the module is
one-dimensional, with \(e=f=0\) and \(h\) an arbitrary scalar
\(\lambda\). If \(\gamma\ne0\), it has a homogeneous basis
\(v,w\) of opposite parities on which
\[
\begin{gathered}
 hv=\lambda v,\qquad hw=(\lambda-1)w,\qquad
 fv=w,\qquad fw=0,\\
 ev=0,\qquad ew=\gamma v,\qquad
 zv=\gamma v,\qquad zw=\gamma w .
\end{gathered}
\]
In each case the scalar data and the parity of \(v\)
determine the isomorphism class in the graded category.
\end{example}

\begin{proof}
The kernel of \(z-\gamma\) for an eigenvalue \(\gamma\)
is a nonzero graded submodule, hence is the whole module.
Since \(e^2=0\), \(\ker e\ne0\). It is graded and
stable under \(h\), so over \(\mathbb C\) it contains a
homogeneous \(h\)-eigenvector \(v\), of weight \(\lambda\).
The table in Proposition~\ref{res:SR02} gives
\(e(fv)=\gamma v\), \(f^2v=0\), and \(h(fv)=(\lambda-1)fv\).
Thus \(\operatorname{span}(v,fv)\) is the whole simple module.
If \(\gamma=0\) and \(fv\ne0\), the line it spans would be a
nonzero proper graded submodule. If \(\gamma\ne0\), the
vector \(w=fv\) is nonzero. The displayed formulas define
a graded representation by the bracket table. It is simple
because either weight line generates both lines, and a nonzero
invariant subspace has a weight vector by interpolation in \(h\).
The weight and parity of the one-dimensional kernel of \(e\)
distinguish the parameters. In this two-dimensional case the
opposite choice of positive odd root instead uses \(\ker f\),
so changes the highest weight to \(\lambda-1\) and changes
its parity. On the one-dimensional modules both choices
give the same highest line.
\end{proof}

\section{The \texorpdfstring{\(p\)}{p}-operation}

For the rest of the chapter, \(k\) has characteristic \(p>0\),
and all Lie algebras are ordinary. Write \(U(\mathfrak g)\)
for the ordinary universal enveloping algebra. We use the
ordered-basis and associated-graded parts of the ordinary PBW
theorem in \emph{Lie Algebras and Finite-Dimensional
Representations}, Result A05~\cite[PDF pp.~8--9]{mg-lie}.
That proof applies to any characteristic for a Lie algebra
with an ordered \(k\)-basis.

\begin{lemma}[Associative \(p\)-powers and the Jacobson polynomials]
\label{res:SR05}
Let \(B\) be an associative \(k\)-algebra, with ordinary
commutator bracket. For \(1\le i\le p-1\), define
\(s_i(x,y)\) by specifying that \(i\,s_i(x,y)\) is the
coefficient of \(t^{i-1}\) in
\((\operatorname{ad}(t x+y))^{p-1}(x)\).
These are Lie polynomials in \(x,y\). One has
\[
 (\operatorname{ad}x)^p=\operatorname{ad}(x^p),\qquad
 (x+y)^p=x^p+y^p+\sum_{i=1}^{p-1}s_i(x,y).
\]
\end{lemma}

\begin{proof}
Left and right multiplication by \(x\) commute.
In characteristic \(p\) the \(p\)-th power of their
difference is the difference of their \(p\)-th powers,
proving the first equality. Set \(X=t x+y\) in \(B[t]\).
Formal differentiation of \(X^p\), using a central variable,
gives
\[
 \frac{d}{dt}X^p=\sum_{j=0}^{p-1}X^{p-1-j}xX^j
          =(\operatorname{ad}X)^{p-1}(x),
\]
since \((-1)^j\binom{p-1}{j}=1\) in \(k\).
The constant and degree-\(p\) coefficients of \(X^p\)
are \(y^p,x^p\). They contribute zero to the derivative.
For each intermediate coefficient, differentiation
multiplies it by the corresponding invertible integer \(i\).
It is therefore \(s_i(x,y)\). Evaluating at \(t=1\)
gives the second equality. The iterated-commutator
description proves that these coefficients are Lie
polynomials and makes their definition available in any
ordinary Lie algebra.
\end{proof}

\begin{definition}[Restricted Lie algebras]
\label{res:SR06}
A restricted Lie algebra is an ordinary Lie algebra with a map
\(x\mapsto x^{[p]}\) satisfying
\[
\begin{aligned}
\operatorname{ad}(x^{[p]})&=(\operatorname{ad}x)^p,\\
(a x)^{[p]}&=a^p x^{[p]},\\
(x+y)^{[p]}&=x^{[p]}+y^{[p]}+\sum_{i=1}^{p-1}s_i(x,y),
\end{aligned}
\]
with the same polynomials as in Lemma~\ref{res:SR05}.
A Lie subalgebra closed under the operation inherits it.
\end{definition}

\begin{proposition}[Matrix and derivation \(p\)-operations]
\label{res:SR07}
Associative \(p\)-th powers restrict to \(\mathfrak{sl}_n(k)\)
inside the full matrix algebra and give a restricted
structure there. They also give a restricted structure on
\(\operatorname{Der}_k(A)\) for a commutative \(k\)-algebra \(A\),
using composition in \(\operatorname{End}_k(A)\).
In particular for \(A=k[t]/(t^p)\) and
\(D_i=t^{i+1}\frac{d}{dt}\), \(-1\le i\le p-2\), one has
\[
 D_0^{[p]}=D_0,\qquad D_i^{[p]}=0\quad(i\ne0).
\]
\end{proposition}

\begin{proof}
Over an algebraic closure a matrix is triangularizable by
inductively choosing an eigenline and passing to the quotient.
Its \(p\)-th power has the \(p\)-th powers of its eigenvalues
on the diagonal. Thus
\(\operatorname{tr}(X^p)=\operatorname{tr}(X)^p\), an identity
which holds already over \(k\). The traceless Lie subalgebra
is therefore closed under the power. For a derivation \(D\),
repeated Leibniz gives
\[
 D^p(ab)=\sum_{j=0}^p\binom pj D^j(a)D^{p-j}(b)
         =D^p(a)b+aD^p(b).
\]
In both cases the restricted identities are those of
Lemma~\ref{res:SR05} and scaling.

Every derivation of \(k[t]/(t^p)\) is determined by the
image of \(t\), which can be arbitrary since the derivative
of \(t^p\) is zero. The displayed \(D_i\) form a basis.
The operator \(D_0\) acts on \(t^m\) by the scalar
\(m\in\mathbb F_p\), whose \(p\)-th power is itself.
On monomials \(D_{-1}\) lowers their degree, and \(D_i\)
with \(i>0\) raises it until the monomial vanishes in
the quotient. Their \(p\)-fold compositions are zero.
For \(p>3\), the simplicity and vanishing Killing form
of this Witt algebra were proved in the earlier volume,
Result J09~\cite[PDF p.~72]{mg-lie}.
\end{proof}

\section{The \texorpdfstring{\(p\)}{p}-center and reduced enveloping algebras}

\begin{theorem}[Polynomial \(p\)-center and reduced PBW]
\label{res:SR08}
Let \(\mathfrak g\) be a restricted Lie algebra of dimension
\(n<\infty\). The function
\[
 z(x)=x^p-x^{[p]}\in U(\mathfrak g)
\]
is central, additive and \(p\)-semilinear:
\(z(a x)=a^p z(x)\). If \(x_1,\ldots,x_n\) is a basis,
the subalgebra \(Z_p\) generated by all \(z(x)\) is the
polynomial ring on \(z(x_1),\ldots,z(x_n)\).
The algebra \(U(\mathfrak g)\) is free over \(Z_p\) on the
ordered monomials with each exponent less than \(p\).

For a linear functional \(\chi:\mathfrak g\to k\), set
\[
 u_\chi(\mathfrak g)
   = U(\mathfrak g)\big/\bigl(z(x)-\chi(x)^p:x\in\mathfrak g\bigr).
\]
The images of those same ordered monomials are a
\(k\)-basis, so the dimension is \(p^n\). In particular
\(u_0(\mathfrak g)\) is the restricted enveloping algebra.
\end{theorem}

\begin{proof}
For \(y\in\mathfrak g\), Lemma~\ref{res:SR05} and the
restricted adjoint identity give
\([x^p,y]=(\operatorname{ad}x)^p(y)=[x^{[p]},y]\);
the \(y\)'s generate \(U(\mathfrak g)\). The two
Jacobson sum formulas for associative powers and the
restricted map cancel upon subtraction, proving additivity.
The scaling identity is immediate. Thus the \(z(x_i)\)
already generate \(Z_p\).

In the associated graded algebra of ordinary PBW,
the leading symbol of \(z(x_i)\) is \(x_i^p\), since
the subtracted term has filtration degree at most one.
Write \(a_i=p b_i+r_i\) with \(0\le r_i<p\) for any
ordinary PBW word of exponents \((a_i)\).
The element
\(\prod_i z(x_i)^{b_i}\prod_i x_i^{r_i}\),
with the \(x_i\) in their fixed order, has that word's
leading symbol. Subtract and use induction on the degree
to express every element as a linear combination of these
products. For independence, distinct pairs \((b,r)\)
have distinct associated-graded leading monomials.
In a finite relation the terms of maximal degree cannot
cancel there. This proves polynomial independence and
the freeness assertion.

By \(p\)-semilinearity the ideal in the definition of
\(u_\chi\) is generated by the central elements
\(z(x_i)-\chi(x_i)^p\). On the free \(Z_p\)-module just
constructed, quotienting by that ideal specializes each
polynomial variable to \(\chi(x_i)^p\). This leaves exactly
the stated basis over \(k\).
\end{proof}

\begin{corollary}[The \(p\)-character of a finite-dimensional simple module]
\label{res:SR09}
If \(k\) is algebraically closed, every finite-dimensional simple
\(\mathfrak g\)-module has a unique linear functional \(\chi\)
such that it is a module over \(u_\chi(\mathfrak g)\).
Conversely a simple module of any \(u_\chi(\mathfrak g)\) is
a finite-dimensional simple \(\mathfrak g\)-module with this
\(p\)-character.
\end{corollary}

\begin{proof}
On a finite-dimensional simple, each central \(z(x_i)\)
acts as a scalar \(c_i\): a nonzero eigenspace is a
submodule. The \(p\)-power map on an algebraically closed
field of characteristic \(p\) is bijective, since it
is injective on a field and every \(T^p-c_i\) has a root.
There is a unique linear \(\chi\) with
\(\chi(x_i)^p=c_i\), and Lemma~\ref{res:SR05} and
Theorem~\ref{res:SR08} ensure the remaining relations
by semilinearity. Conversely \(u_\chi\) is finite-dimensional
and a simple module is generated by any one nonzero vector,
so is finite-dimensional. Submodules for the quotient are
the same as invariant subspaces for the Lie generators,
proving simplicity and the assertion about \(\chi\).
\end{proof}

For the restricted identities and reduced algebra see
Jantzen~\cite[Sections A.1--A.3 and A.7, PDF pp.~1--3 and 4--5]{jantzen-modular}
and Strade--Farnsteiner~\cite[Sections 2.1, 2.5 and 5.3]{stradefarnsteiner}.

\section{Baby Verma modules for \texorpdfstring{\(\mathfrak{sl}_2\)}{sl(2)}}

Assume now that \(k\) is algebraically closed and \(p>2\).
Put \(e=E_{12}\), \(f=E_{21}\), and \(h=E_{11}-E_{22}\).
The brackets are \([h,e]=2e\), \([h,f]=-2f\), and
\([e,f]=h\). By Proposition~\ref{res:SR07} one has
\(e^{[p]}=f^{[p]}=0\) and \(h^{[p]}=h\).
This section computes the simple reduced modules in all
three types of \(p\)-character, separating their highest
weight labels from the characteristic-zero labels.

\begin{proposition}[Baby Verma induction]
\label{res:SR10}
Suppose \(\chi(e)=0\). Write \(\eta=\chi(h)\) and
\(\xi=\chi(f)\). For each root \(\lambda\) of
\(\lambda^p-\lambda=\eta^p\), there is a module \(Z_\chi(\lambda)\)
with basis \(v_i=f^i v_0\), \(0\le i<p\), and
\[
\begin{aligned}
 hv_i&=(\lambda-2i)v_i,\\
 ev_0&=0,\qquad ev_i=i(\lambda-i+1)v_{i-1}\quad(1\le i<p),\\
 fv_i&=v_{i+1}\quad(i<p-1),\qquad fv_{p-1}=\xi^p v_0.
\end{aligned}
\]
Every simple \(u_\chi(\mathfrak{sl}_2)\)-module is a
quotient of one of these modules.
\end{proposition}

\begin{proof}
For \(\mathfrak b=k h+k e\), the one-dimensional space
with \(h\) acting by \(\lambda\) and \(e\) by zero respects
the bracket and the reduced power identities, so is a
\(u_{\chi|_{\mathfrak b}}(\mathfrak b)\)-module. Here
\(\mathfrak b\) is closed under the \(p\)-map because powers
of an upper triangular matrix stay upper triangular.
Order the reduced PBW monomials as \(f^i h^j e^\ell\).
The basis in Theorem~\ref{res:SR08} identifies the reduced
algebra of \(\mathfrak b\) with its subalgebra in the
reduced algebra of \(\mathfrak{sl}_2\), and shows the
latter is free as a right module over it on \(f^i\).
Inducing the one-dimensional module therefore gives a
\(p\)-dimensional module with the asserted basis.
The \(h\) and \(f\) formulas follow from the bracket
and the defining reduced relations. Commuting \(e\)
past one further \(f\), the coefficient in
\(e f^{i+1}v_0\) is
\((\lambda-2i)+i(\lambda-i+1)=(i+1)(\lambda-i)\).
This proves the \(e\) formula by induction.

On any simple \(u_\chi(\mathfrak{sl}_2)\)-module \(e^p=0\), so
\(\ker e\ne0\), and \([h,e]=2e\) implies this kernel
is \(h\)-stable. Choose an \(h\)-eigenvector in it.
Its weight satisfies the reduced relation for \(h\).
The induced module with that weight maps onto the simple,
because the nonzero vector generates it.
\end{proof}

\begin{theorem}[The zero \(p\)-character]
\label{res:SR11}
For \(\chi=0\) the simple modules are \(L(m)\) for
\(m=0,\ldots,p-1\). They have dimensions \(m+1\).
For \(\lambda=m\) in the prime field, \(L(m)\) is the
quotient of \(Z_0(m)\) by the span of
\(v_{m+1},\ldots,v_{p-1}\), with empty span for \(m=p-1\).
\end{theorem}

\begin{proof}
The solutions of \(\lambda^p-\lambda=0\) are the prime
field elements. The \(p\) weights of \(Z_0(m)\) are
distinct because \(2\ne0\) in \(k\). The displayed tail
is stable: its first vector is killed by \(e\), and \(f\)
kills \(v_{p-1}\). In the quotient a nonzero submodule
contains a weight vector \(v_i\) with \(0\le i\le m\),
since projections onto the finitely many weight spaces
are polynomials in \(h\). For this range every noninitial
\(e\)-coefficient is nonzero. Repeated application of \(e\)
reaches \(v_0\) and \(f\) then produces the entire quotient.
This proves simplicity. The same argument in \(Z_0(m)\)
shows that any submodule containing a vector with nonzero
component in one of the first \(m+1\) weight spaces is
the whole module. Thus its proper submodules all lie in
the tail and the simple quotient is unique.
The modules \(L(m)\) are pairwise nonisomorphic by their
distinct dimensions, and Proposition~\ref{res:SR10}
gives exhaustion.
\end{proof}

\begin{theorem}[Nonzero semisimple and nilpotent \(p\)-characters]
\label{res:SR12}
There are the following two nonzero cases, after a change
of basis by conjugation of matrices.

\emph{Semisimple case.} Suppose
\(\chi(e)=\chi(f)=0\), \(\eta=\chi(h)\ne0\).
The \(p\) modules \(Z_\chi(\lambda)\) for the roots of
\(\lambda^p-\lambda=\eta^p\) are simple and pairwise
nonisomorphic, each of dimension \(p\).

\emph{Nilpotent case.} Suppose
\(\chi(e)=\chi(h)=0\), \(\chi(f)=1\).
All \(p\) baby Verma modules \(Z_\chi(m)\), \(m\in\mathbb F_p\),
are simple of dimension \(p\). The isomorphism classes
are indexed by the orbits of \(m\mapsto -m-2\), giving
\((p+1)/2\) classes.
\end{theorem}

\begin{proof}
In the semisimple case the polynomial of the weights is
separable of degree \(p\) and has no root in the prime field.
All \(e\)-coefficients for \(1\le i<p\) in
Proposition~\ref{res:SR10} are nonzero. A nonzero
submodule contains a weight line, and repeated applications
of \(e\) and \(f\) generate the entire basis. The kernel of
\(e\) is only the highest line, whose weight distinguishes
the \(p\) modules. They exhaust by induction.

In the nilpotent case the \(h\)-weights are still distinct
and \(f\) cyclically permutes the \(p\) weight lines since
\(f^p=1\). A submodule containing one weight line therefore
contains them all. For \(m\ne-1\) the kernel of \(e\) is
spanned by \(v_0,v_{m+1}\), where \(m\) is represented by
its integer in \(0,\ldots,p-2\). Their weights are \(m\)
and \(m-2(m+1)=-m-2\). Each is a highest line and induces
a map onto the simple, so the modules with these two labels
are isomorphic. For \(m=-1\) only \(v_0\) is killed by \(e\).
Every isomorphism preserves the weights in \(\ker e\), so
there are no other identifications. There is exactly one
fixed element of the involution, \(-1\), giving the count.

Finally the trace pairing on \(\mathfrak{sl}_2(k)\) is
nondegenerate: it pairs \(e\) with \(f\) and has
\(\operatorname{tr}(h^2)=2\). Every \(\chi\) is therefore
\(\operatorname{tr}(X\,\cdot)\) for one traceless \(X\).
A nonzero such matrix is conjugate either to
\(\operatorname{diag}(a,-a)\) with \(a\ne0\) or to
\(E_{12}\), according as its two eigenvalues are distinct
or both zero. The displayed parameters follow from the
pairing. Conjugation preserves the bracket and matrix
\(p\)-power, so it transports the corresponding reduced
algebras and their modules. This proves the scope of the
two cases.
\end{proof}

The inductive construction follows the setting of
Jantzen~\cite[Sections B.3--B.4, B.10 and C.3,
PDF pp.~6--7 and 9--11]{jantzen-modular}.
General classification of simple modular Lie algebras and of
simple modules for higher-rank reduced enveloping algebras
uses further hypotheses and results. Here the module
classifications are for \(\mathfrak{gl}(1|1)\) over
\(\mathbb C\) and for ordinary \(\mathfrak{sl}_2\) over an
algebraically closed field of characteristic greater than two.

\chapter[The Type A2 Quantum Enveloping Algebra]{The Type \(A_2\) Quantum Enveloping Algebra}
\label{ch:quantum-enveloping}

We construct the two-root quantum enveloping algebra and its formal
classical limit. The transition from one simple root to two requires
the quantum Serre relations as well as the Cartan and mixed
commutator relations. The rank-one algebra, its Hopf structure,
and its finite-dimensional type-1 representations are treated
in Results H10--H15 of \emph{Topics in Rings, Modules, and
Associative Algebras}~\cite[PDF pp.~25--28]{mg-rm03d}.
Here the main tasks are to prove a basis for the two-root algebra,
check the coproduct on the new relations, and compute its
first-order skew part. The constructions below use
\(\Delta^{\mathrm{op}}\) for tensor factor exchange applied to
\(\Delta\), and \(a\wedge b=a\otimes b-b\otimes a\).

\section{The quantum Serre relations and the positive part}

Let \(A\) be a commutative coefficient ring and \(q\in A^\times\).
Write \(c=q+q^{-1}\) and put
\[
 \mathcal S(x,y)=x^2y-cxyx+yx^2 .
\]

\begin{theorem}[The two-generator positive part]
\label{res:UQ01}
The \(A\)-algebra on \(x,y\) with relations
\(\mathcal S(x,y)=\mathcal S(y,x)=0\) is free as an
\(A\)-module with basis \(x^a z^b y^d\), where
\[
 z=xy-q^{-1}yx,\qquad a,b,d\in\mathbb Z_{\ge0}.
\]
In this algebra \(xz=qzx\), \(zy=qyz\), and \(yx=qxy-qz\).
This basis and presentation commute with arbitrary base change
on \(A\) carrying \(q\) to a unit.
\end{theorem}
\begin{proof}
The expressions \(xz-qzx\) and \(zy-qyz\) expand to
\(\mathcal S(x,y)\) and \(\mathcal S(y,x)\), respectively.
The remaining relation is the defining formula for \(z\)
solved for \(yx\). Conversely these three relations for
\(x,z,y\) eliminate \(z\) and give precisely the two Serre
relations.

Adjoin \(z\) to \(A[x]\) with \(zx=q^{-1}xz\).
On the resulting algebra let \(\sigma(x)=qx\) and
\(\sigma(z)=q^{-1}z\). This preserves the relation and
has inverse given by reciprocal scalars. Define the
\(\sigma\)-derivation by \(\delta(x)=-qz\), \(\delta(z)=0\),
where \(\delta(ab)=\sigma(a)\delta(b)+\delta(a)b\).
It respects the relation because both \(\delta(zx)\) and
\(\delta(q^{-1}xz)\) are \(-z^2\). Hence adjoining \(y\)
with \(ya=\sigma(a)y+\delta(a)\) gives the second pair
of commutation rules. The Ore normal form in
\emph{Topics in Rings, Modules, and Associative Algebras},
Result O11~\cite[PDF p.~34]{mg-rm03d}, makes the indicated
monomials a free basis. Every relation and each basis
applies over the new coefficient ring by the same construction.
\end{proof}

At \(q=1\), the generator \(z\) is \([x,y]\) and commutes with
\(x,y\). In the matrix realization of \(\mathfrak{sl}_3\), the
assignment \(x\mapsto e_{12}\), \(y\mapsto e_{23}\),
\(z\mapsto e_{13}\) then identifies this presentation
with the enveloping algebra of the positive nilpotent subalgebra.
The commutation rules and the basis on each side prove the
identification, using the ordered PBW basis of
\emph{Lie Algebras and Finite-Dimensional Representations},
Result A05~\cite[PDF pp.~8--9]{mg-lie}, on the classical side.

\section{The Cartan coefficients and triangular normal form}

The common Cartan matrix will be \(C=(a_{ij})_{1\le i,j\le 2}\),
with \(a_{11}=a_{22}=2\) and \(a_{12}=a_{21}=-1\).
We give parallel coefficient choices so that the normal-form
argument can also be used to justify the formal limit.

\begin{definition}[The generic and truncated formal presentations]
\label{res:UQ02}
Fix a field \(k\) of characteristic zero. The generic choice is
\(A=k\) with a parameter \(q\in k^\times\), \(q^2\ne1\), or
\(A=k(q)\) with \(q\) indeterminate. It uses
\[
 T=A[K_1^{\pm1},K_2^{\pm1}],\qquad
 \tau_i(K_l)=q^{a_{li}}K_l,\qquad
 D_i=\frac{K_i-K_i^{-1}}{q-q^{-1}} .
\]
The other choice, for an integer \(N\ge1\), uses
\[
 \begin{gathered}
 A=k[h]/(h^N),\qquad T=A[H_1,H_2],\qquad
 q=e^h,\qquad K_i=e^{hH_i},\\
 \tau_i(H_l)=H_l+a_{li},\qquad
 D_i=\frac{\sinh(hH_i)/h}{\sinh(h)/h}\pmod {h^N}.
 \end{gathered}
\]
The divided series are first formed in \(k[H_1,H_2][[h]]\);
the denominator there has constant term \(1\).
Both choices have \(\tau_i(K_l)=q^{a_{li}}K_l\), and the
\(\tau_i\) commute. The basic relations on \(T\) and \(E_i,F_i\) are
\[
 t F_j=F_j\tau_j^{-1}(t),\qquad
 E_i t=\tau_i^{-1}(t)E_i,\qquad
 [E_i,F_j]=\delta_{ij}D_i \quad (t\in T).
\]
The quantum enveloping algebra is obtained by also imposing
\(\mathcal S(E_i,E_j)=\mathcal S(F_i,F_j)=0\) for \(i\ne j\).
Write \(U\) for the generic choice and \(U_N\) for the truncated
formal choice.
\end{definition}

\begin{lemma}[Normal form before the Serre quotient]
\label{res:UQ03}
With the coefficients of Definition~\ref{res:UQ02},
let \(B_0\) impose only the basic relations on \(T,E_i,F_i\).
Multiplication identifies \(B_0\), as an \(A\)-module, with
\[
 A\langle F_1,F_2\rangle\otimes_A T
       \otimes_A A\langle E_1,E_2\rangle .
\]
There is an algebra involution \(\omega\) interchanging
\(E_i,F_i\) and inverting the Cartan generators
(\(K_i\mapsto K_i^{-1}\) or \(H_i\mapsto-H_i\)).
There is also an algebra anti-involution \(\eta\) interchanging
\(E_i,F_i\) and fixing \(T\).
\end{lemma}
\begin{proof}
Use the basic relations to move all negative letters left of
the \(T\)-coefficient and all positive letters right of it,
and combine adjacent \(T\)-coefficients. To see termination,
first count the total number of \(E,F\) letters, then the
total number of factors, and then inversions relative to the
order negative, Cartan, positive. Replacing \(E_iF_j\) by its
commutator term lowers the first count. Combining adjacent
Cartan factors lowers the second. Every remaining interchange
lowers the inversion count. One can start with Cartan factors
from a monomial basis of \(T\) and expand all replacements
linearly.

Disjoint replacements commute. Overlaps with adjacent
Cartan factors are compatible with multiplication since the
\(\tau_i\) are algebra automorphisms. The only remaining
overlap is \(E_i tF_j\). In the two orders the leading term
is \(F_j\tau_i^{-1}\tau_j^{-1}(t)E_i\); the other term in
one order is \(\delta_{ij}\tau_i^{-1}(t)D_i\) and in
the other is \(\delta_{ij}D_i\tau_j^{-1}(t)\).
These coincide by commutativity of the translations and
of \(T\). The normal form is therefore independent of reductions:
induct on the ordered counts and compare the first two choices,
using the overlap computation or commuting them when disjoint.
Normalizing all terms in a defining relation inserted into
a word gives zero, so the normal form descends to the
presented algebra and proves independence as well as spanning.

Cartan inversion \(\kappa\) satisfies
\(\kappa\tau_i=\tau_i^{-1}\kappa\) and \(\kappa(D_i)=-D_i\).
These identities check the two Cartan and mixed relations
after the swap defining \(\omega\). Reversing all products
while interchanging \(E,F\) gives the Cartan relations in
the opposite order, and sends
\([E_i,F_j]\) to \([E_j,F_i]\). This checks \(\eta\).
\end{proof}

\begin{lemma}[Crossing the Serre ideal]
\label{res:UQ04}
In the algebra \(B_0\) of Lemma~\ref{res:UQ03},
\(\mathcal S(E_i,E_j)\) commutes with every \(F_l\), and
\(\mathcal S(F_i,F_j)\) commutes with every \(E_l\).
Under the same normal form, the two-sided ideal generated
by the positive and negative Serre polynomials is the sum
of the submodules obtained by putting the corresponding
free-algebra Serre ideal in a positive or negative factor.
\end{lemma}
\begin{proof}
Put \(S=\mathcal S(E_i,E_j)\), \(i\ne j\).
For \(l=j\) use
\(
 [S,F_j]=(\tau_i^{-2}D_j-c\tau_i^{-1}D_j+D_j)E_i^2.
\)
For \(l=i\), the commutator rule
\([ab,F]=a[b,F]+[a,F]b\) gives
\[
\begin{split}
 [S,F_i]
 &=\bigl(D_i+\tau_i^{-1}D_i-c\tau_i^{-1}\tau_j^{-1}D_i\bigr)E_iE_j\\
 &\quad+\bigl(\tau_j^{-1}(D_i+\tau_i^{-1}D_i)-cD_i\bigr)E_jE_i .
\end{split}
\]
Every displayed coefficient is zero. In the generic case,
clear \(q-q^{-1}\) and evaluate each coefficient on
the terms \(K_l^s\) with \(s=1,-1\).
The three scalar identities are
\[
 q^{2s}-c q^s+1=0,\qquad
 1+q^{-2s}-c q^{-s}=0,\qquad
 q^s(1+q^{-2s})-c=0 .
\]
They follow directly from \(c=q+q^{-1}\).
In the formal case do the same calculation on
\(K_l^s=e^{shH_l}\) before reduction.
The relations with \(D_l\) then follow by cancellation
of \(e^h-e^{-h}\) in the formal coefficient domain,
so they hold in every truncation. Applying \(\omega\)
gives the negative commutators.

To check the ideal statement, take the submodule with
positive Serre ideal in the positive factor. Left
multiplication by \(F_l\) and right multiplication by \(E_l\)
preserve it at once. Cartan factors shift across a word
according to the total number of its root letters; a Serre
polynomial is homogeneous for those counts.
Thus left or right Cartan multiplication preserves
the submodule. Moving a left \(E_l\) across a negative
word uses only the basic relations and produces either
a remaining positive letter or shorter negative words
with Cartan coefficients. This also preserves the positive
ideal factor. For right multiplication by \(F_l\), it
suffices to move it across a positive expression \(pSq\).
The commutators \([p,F_l]\) and \([q,F_l]\) have only
Cartan coefficients and positive words. The commutator
\([S,F_l]\) is zero, and moving a Cartan coefficient
across \(S\) preserves its Serre factor by homogeneity.
This proves closure on both sides. Applying \(\eta\) gives
the same result for the negative submodule without changing
the factor order. Their sum is a two-sided ideal containing
the generators, and all its basis-type terms already
belong to the generated ideal.
\end{proof}

\begin{theorem}[Triangular PBW in type \(A_2\)]
\label{res:UQ05}
For either coefficient choice of Definition~\ref{res:UQ02},
multiplication is an \(A\)-module isomorphism from
the negative Serre algebra tensored with \(T\) and then
with the positive Serre algebra onto the quantum
enveloping algebra. In particular, writing
\(z_-=F_1F_2-q^{-1}F_2F_1\) and
\(z_+=E_1E_2-q^{-1}E_2E_1\), a basis consists of
\[
 F_1^a z_-^b F_2^d\, T_0\, E_1^{a'} z_+^{b'} E_2^{d'},
\]
where the six root exponents are nonnegative integers and
\(T_0\) ranges over the indicated Cartan monomials.
In the generic case the subalgebra generated by
\(F_i,K_i^{\pm1},E_i\) is the rank-one algebra
\(U_q(\mathfrak{sl}_2)\).
\end{theorem}
\begin{proof}
Apply Lemmas~\ref{res:UQ03} and \ref{res:UQ04};
passing to the Serre quotient in either factor gives
the asserted tensor product by right exactness of
tensoring. Theorem~\ref{res:UQ01} supplies the bases
of the two Serre factors, and \(T\) has its evident
monomial basis. The fixed-index generators satisfy the
rank-one relations, and the corresponding monomials
\(F_i^aK_i^bE_i^d\) are a subset of the independent
displayed family. This proves that the rank-one map is
injective by the rank-one PBW result H10
~\cite[PDF p.~25]{mg-rm03d}.
\end{proof}

\section{The Hopf structure}

\begin{theorem}[Coproduct and antipode in type \(A_2\)]
\label{res:UQ06}
The formulas
\[
\begin{aligned}
 \Delta(E_i)&=E_i\otimes1+K_i\otimes E_i, &
 \Delta(F_i)&=F_i\otimes K_i^{-1}+1\otimes F_i,\\
 \epsilon(E_i)&=\epsilon(F_i)=0, &
 S(E_i)&=-K_i^{-1}E_i,\qquad S(F_i)=-F_iK_i
\end{aligned}
\]
give Hopf structures on \(U\) and \(U_N\) from
Definition~\ref{res:UQ02}. In the generic case \(K_i\)
is group-like with counit \(1\) and antipode \(K_i^{-1}\).
In the formal case \(H_i\) is primitive with counit
zero and antipode \(-H_i\), so \(K_i=e^{hH_i}\) is
group-like. The antipodes are bijective, and the formal
structures commute with reduction in \(N\).
\end{theorem}
\begin{proof}
The proposed Cartan coproduct preserves its own relations.
For the generic case, conjugation by \(K_l\otimes K_l\)
scales the images of \(E_i,F_i\) by \(q^{a_{li}},q^{-a_{li}}\).
For the formal case, commuting with \(H_l\otimes1+1\otimes H_l\)
gives the same two weights in additive form. Thus the
Cartan-root relations hold. In the commutator of
\(\Delta(E_i),\Delta(F_j)\) the mixed term
\([K_i\otimes E_i,F_j\otimes K_j^{-1}]\) vanishes:
the factors needed to interchange the two products
are \(q^{-a_{ij}}\) and \(q^{a_{ji}}\), which cancel.
The term \([E_i\otimes1,1\otimes F_j]\) also vanishes.
The other terms are therefore
\[
 \delta_{ij}\bigl(D_i\otimes K_i^{-1}+K_i\otimes D_i\bigr)
       =\delta_{ij}\Delta(D_i).
\]
For Laurent polynomials this last identity follows from
multiplying by \(q-q^{-1}\) and using \(\Delta(K_i)=K_i\otimes K_i\).
In the formal case it follows by the same calculation
before cancellation of \(e^h-e^{-h}\) in the
formal series ring in the two Cartan sets, then reduction.

For the positive Serre relation take \(i\ne j\) and put
\[
 x=E_i\otimes1,\quad u=K_i\otimes E_i,\qquad
 y=E_j\otimes1,\quad v=K_j\otimes E_j.
\]
We need \(\mathcal S(x+u,y+v)=0\). The pure parts are
\(\mathcal S(x,y)\) and \(\mathcal S(u,v)\); the latter
is zero because the first-factor \(K\)'s commute and
the second factor is the Serre polynomial.
Use \(ux=q^2xu,\ uy=q^{-1}yu,\ vx=q^{-1}xv\).
Terms with \(y\) and two copies of \(u\), or with \(v\) and two
copies of \(x\), have scalar factor \(1-cq^{-1}+q^{-2}=0\).
Terms with \(y\) and one copy each of \(x,u\) sum to
\[
 cxyu-c(xyu+qyxu)+(1+q^2)yxu=0.
\]
Those with \(v\) and one copy each of \(x,u\) sum to
\[
 (1+q^2)xuv-c(xvu+qxuv)+cxvu=0 .
\]
This covers every term. The negative relation follows
by \(\omega\) from Lemma~\ref{res:UQ03} on both factors
and then tensor-factor exchange. Those transformations
send the displayed positive coproduct formula to the
negative one. Hence \(\Delta\) is an algebra map.

The counit respects the relations, including \(\epsilon(D_i)=0\).
Coassociativity follows on the generators: for \(E_i\) each
iterate is \(E_i\otimes1\otimes1+K_i\otimes E_i\otimes1+
K_i\otimes K_i\otimes E_i\); the \(F_i\) expansion similarly
puts \(K_i^{-1}\) after each \(F_i\). The formulas for group-like
and primitive elements are coassociative directly. The counit
identities also hold on the generators.

To check that \(S\) extends anti-multiplicatively, use the
Cartan inversion \(\kappa\) of Lemma~\ref{res:UQ03}.
The identity \(\kappa\tau_i=\tau_i^{-1}\kappa\) checks
the reversed Cartan-root relations, since \(S(E_i)\)
and \(S(F_i)\) have their respective original weights.
For the mixed relation, move \(K_i^{-1}\) to the left
and \(K_j\) to the right. Symmetry of \(a_{ij}\) gives
\[
 [S(F_j),S(E_i)]
    =K_i^{-1}(F_jE_i-E_iF_j)K_j=-\delta_{ij}D_i,
\]
which is the required value of \(S(\delta_{ij}D_i)\).
For each reversed word in a positive Serre polynomial,
move all \(K^{-1}\)'s to the left. The scalar in addition
to the three minus signs is the product of \(q^{a_{rs}}\)
over the three pairs of letters. This is independent of
their order because \(a_{rs}=a_{sr}\). Here the pair sum
is \(2-1-1=0\). The result is the common Cartan monomial
times the reversed Serre polynomial, which is zero.
The negative check moves all \(K\)'s to the right,
with the reciprocal scalar, and gives the same conclusion.

The two antipode products give the counit times \(1\)
on every generator by substitution. They extend to products
by anti-multiplicativity: for instance in
\(\sum S((ab)_{(1)})(ab)_{(2)}\) the two contributions of
\(a\) occur next to one another, giving \(\epsilon(a)\),
and then the contributions of \(b\) give \(\epsilon(b)\).
For the other product the contributions of \(b\) occur
first. Finally \(S^2\) fixes \(T\) and scales \(E_i,F_i\)
by \(q^{-2},q^2\), respectively. These reciprocal scalings
preserve the defining relations and have inverses.
Thus \(S^2\) is invertible and \(S\) is bijective.
All formulas in the formal choice use reductions of the
same series, so they commute with reduction in \(N\).
\end{proof}

\section{The formal classical limit}

\begin{theorem}[A topologically free quantization of the enveloping algebra]
\label{res:UQ07}
Use the compatible truncations of Definition~\ref{res:UQ02}
and set \(U_h=\varprojlim_N U_N\).
It is a topologically free \(k[[h]]\)-module carrying the
completed Hopf structure of Theorem~\ref{res:UQ06}.
Its reduction modulo \(h\) is \(U(\mathfrak{sl}_3(k))\),
as a Hopf algebra. Here the completed tensor product
is the inverse limit of \(U_N\otimes_{k[h]/(h^N)} U_N\).
\end{theorem}
\begin{proof}
The monomials in Theorem~\ref{res:UQ05} have compatible
reductions. Thus each \(U_N\) is the free module on the
same index set over \(k[h]/(h^N)\), and reduction is
coefficientwise. An element of the limit is a series
\(\sum_{n\ge0}h^n v_n\), where each \(v_n\) is a finite
\(k\)-linear combination of the indexed monomials.
Indeed, a finite truncation has finite support and
conversely a finite list of such \(v_n\)'s has finite
total support. This proves topological freeness and
identifies the quotient by \(h\) with \(U_1\).
The same statement with pairs of basis indices proves
topological freeness of the completed tensor product.
The Hopf maps are the compatible maps in the theorem
cited above, so their identities pass to inverse limits.

Modulo \(h\), the Cartan commutators in Definition~\ref{res:UQ02}
give \([H_l,E_i]=a_{li}E_i\), \([H_l,F_i]=-a_{li}F_i\),
and \([E_i,F_j]=\delta_{ij}H_i\). The Serre relations
become the ordinary double commutator relations.
The assignment to \(\mathfrak{sl}_3(k)\) sends
\[
 E_1\mapsto e_{12},\quad E_2\mapsto e_{23},\quad
 F_1\mapsto e_{21},\quad F_2\mapsto e_{32},\quad
 H_i\mapsto e_{ii}-e_{i+1,i+1},
\]
where \(e_{ij}\) is a matrix unit. These matrices satisfy
the displayed relations by matrix multiplication.
The resulting algebra map to \(U(\mathfrak{sl}_3(k))\)
sends \(z_+,z_-\) to \(e_{13},-e_{31}\).
Its basis maps precisely to the ordered PBW monomials
for this basis of \(\mathfrak{sl}_3(k)\), by Result A05 of
\emph{Lie Algebras and Finite-Dimensional Representations}
~\cite[PDF pp.~8--9]{mg-lie}. Thus the map is an
isomorphism. The coproduct reduces to the primitive formula
on the indicated generators, and the counit and antipode
reduce to the standard ones. This also proves the
Hopf assertion.
\end{proof}

\begin{theorem}[The first-order cobracket on \(\mathfrak{sl}_3\)]
\label{res:UQ08}
On the specialization of Theorem~\ref{res:UQ07}, use
lowercase \(e_i,f_i,h_i\) for the images of \(E_i,F_i,H_i\).
For a lift \(X\in U_h\) of \(x\), put
\[
 \delta(x)=\frac{\Delta(X)-\Delta^{\mathrm{op}}(X)}h
             \pmod h .
\]
Its restriction to \(\mathfrak{sl}_3(k)\) is a Lie
bialgebra cobracket. We identify \(x\wedge y\) with
\(x\otimes y-y\otimes x\). Explicitly,
\[
\begin{aligned}
 \delta(h_i)&=0, &
 \delta(e_i)&=h_i\wedge e_i, &
 \delta(f_i)&=h_i\wedge f_i,\\
 \delta([e_1,e_2])&=(h_1+h_2)\wedge[e_1,e_2]+2e_1\wedge e_2,\\
 \delta([f_1,f_2])&=(h_1+h_2)\wedge[f_1,f_2]-2f_1\wedge f_2 .
\end{aligned}
\]
\end{theorem}
\begin{proof}
The specialization is cocommutative, so the numerator
is divisible by \(h\). The completed tensor product
is \(h\)-torsion-free by Theorem~\ref{res:UQ07}.
Replacing a lift by \(X+hY\) does not change the
answer, since the specialization of
\(\Delta(Y)-\Delta^{\mathrm{op}}(Y)\) is zero.
Expanding \(K_i=1+hH_i+O(h^2)\) and
\(K_i^{-1}=1-hH_i+O(h^2)\) in Theorem~\ref{res:UQ06}
gives the displayed simple-root and Cartan values.

Write \(\Delta_0\) for the specialized coproduct.
Multiplicativity of \(\Delta,\Delta^{\mathrm{op}}\)
and division of their skew difference give
\[
 \delta([x,y])=[\Delta_0(x),\delta(y)]
                  +[\delta(x),\Delta_0(y)] .
\]
For primitive \(x,y\) the right-hand side is
\(\operatorname{ad}_x\delta(y)-\operatorname{ad}_y\delta(x)\)
for the adjoint action on both factors. Using the
two off-diagonal Cartan entries \(-1\) gives
the two remaining displayed formulas. The eight
matrix basis elements from the preceding proof are thus
sent into the skew part of \(\mathfrak{sl}_3\otimes\mathfrak{sl}_3\).
The commutator identity gives the Lie-algebra cocycle rule
on that subspace.

Here is the co-Jacobi check. Write \(B=\Delta-\Delta^{\mathrm{op}}\)
and \(D=(\Delta\otimes1)\Delta=(1\otimes\Delta)\Delta\).
In tensor notation, write \(D\) as \(123\) and denote
its factor permutations by their orders. Expansion gives
\((B\otimes1)B=123-213-231+321\).
The sum of its three cyclic permutes is zero, because
\(123,231\) are cyclic permutations and \(213,321\) are
also cyclic permutations. The same coefficientwise
description of the triple completed tensor product
shows it is \(h\)-torsion-free. Dividing by \(h^2\)
and reducing gives the co-Jacobi identity for \(\delta\).
Pairing this cyclic identity with triples of linear forms
gives the Jacobi identity for the dual bracket. Together
with the skew and cocycle rules this proves the Lie
bialgebra assertion.
\end{proof}

\begin{corollary}[The coboundary tensor in the matrix model]
\label{res:UQ09}
With the convention of Theorem~\ref{res:UQ08},
the same cobracket on \(\mathfrak{sl}_3\) is
\(\delta(x)=\operatorname{ad}_x(r)\), where
\[
 r=-e_{12}\wedge e_{21}-e_{23}\wedge e_{32}-e_{13}\wedge e_{31}.
\]
\end{corollary}
\begin{proof}
The matrix-unit rule is
\([e_{ij},e_{lm}]=\delta_{jl}e_{im}-\delta_{mi}e_{lj}\).
Each wedge in \(r\) has total weight zero, giving
\(\operatorname{ad}_{h_i}(r)=0\).
For \(e_1=e_{12}\), the first wedge with its minus sign
contributes \(h_1\wedge e_1\); the other two contributions
are \(-e_{13}\wedge e_{32}\) and \(e_{13}\wedge e_{32}\).
The same computation for \(e_2=e_{23}\) gives
\(h_2\wedge e_2\), with cross contributions
\(e_{13}\wedge e_{21}\) and \(-e_{13}\wedge e_{21}\).
The Lie algebra automorphism \(x\mapsto-x^{\mathsf T}\)
sends \(r\) to \(-r\), \(e_i\) to \(-f_i\) and
\(h_i\) to \(-h_i\). Applying it to the identities for
\(e_i\) gives \(\operatorname{ad}_{f_i}(r)=h_i\wedge f_i\).
Both cobrackets satisfy the cocycle identity, for
\(\operatorname{ad}_x(r)\) by the Jacobi identity.
Their equality on the six generators therefore
implies equality on all of \(\mathfrak{sl}_3\).
\end{proof}

\paragraph{Sources.}
Kassel~\cite[Sections XVI.5 and XVII.2]{kassel} gives the
general Drinfeld--Jimbo constructions. Etingof--Semenyakin
~\cite[Sections 4.1--4.2]{etingof-semenyakin} treats the
Serre presentation and quantization in general.
This chapter proves the full normal form and formal limit
for type \(A_2\) directly, with a specified skew-coproduct
convention.

\chapter{Cohomology and Explicit Deformations}

Work over a field \(k\) of characteristic zero. For Lie brackets,
our convention for the Chevalley--Eilenberg differential, changes
of basis, and obstruction classes is that of
\emph{Lie Algebras and Finite-Dimensional Representations},
Results I02--I03~\cite[PDF pp.~56--57]{mg-lie}.
We use those results to compute an entire first-order obstruction
locus in a small example. We then turn to Lie bialgebras and
deformations of the coproduct. In that setting multiplication
and comultiplication must satisfy a compatibility equation
in addition to their separate associativity conditions.

\section{An adjoint cohomology calculation}

\begin{proposition}[Adjoint cohomology of the three-dimensional Heisenberg algebra]
\label{res:LD01}
Let \(\mathfrak h\) have basis \(x,y,z\) with
\([x,y]=z\) and \(z\) central. Then
\(\dim H^2(\mathfrak h,\mathfrak h)=5\) and
\(\dim H^3(\mathfrak h,\mathfrak h)=2\), with adjoint
coefficients. Each class in degree two has a unique representative
\(\phi\) of the form
\[
 \phi(x,y)=0,\qquad
 \phi(x,z)=a x+b y+c z,\qquad
 \phi(y,z)=d x-a y+e z .
\]
In degree three evaluation on \((x,y,z)\) identifies
\(H^3\) with \(\mathfrak h/kz\).
\end{proposition}
\begin{proof}
For an alternating \(2\)-cochain \(\psi\), write
\(A=\psi(x,y)\), \(B=\psi(x,z)\), and \(C=\psi(y,z)\).
The argument with a bracket in the Chevalley--Eilenberg
formula is zero: the only nonzero bracket of two basis vectors
in the triple is \([x,y]=z\), and \(\psi(z,z)=0\).
Hence
\[
 (d\psi)(x,y,z)=[x,C]-[y,B]
               =(C_y+B_x)z .
\]
Here subscripts denote coordinates in the fixed basis.
Thus the cocycle space has dimension \(8\). The image in
degree three is exactly \(kz\), and there are no degree-four
cochains, proving the degree-three assertion.

For a linear \(1\)-cochain \(f\) write
\(f(z)=u x+v y+w z\). Its differential is
\[
 \begin{aligned}
 (df)(x,y)&=-u x-v y+(f(y)_y+f(x)_x-w)z,\\
 (df)(x,z)&=v z,\\
 (df)(y,z)&=-u z.
 \end{aligned}
\]
Thus a coboundary is determined by its value on \((x,y)\),
and that value ranges over all of \(\mathfrak h\), for instance
by taking \(f(x)=f(y)=0\) and varying \(f(z)\).
Subtract the unique coboundary matching \(A\). In the resulting
normalized cocycle \(A=0\), the only equation is
\(C_y=-B_x\). This gives the claimed parameters and uniqueness.
\end{proof}

\begin{theorem}[Exactly which of these first-order classes extend]
\label{res:LD02}
In the coordinates of Proposition~\ref{res:LD01}, a first-order
bracket deformation of \(\mathfrak h\) extends modulo \(t^3\)
if and only if
\[
                  cd-ea=0,\qquad eb+ca=0 .
\]
When both equations hold, the bracket \([\, ,\,]+t\phi\)
already satisfies Jacobi over \(k[t]\), and in particular gives
a formal extension over \(k[[t]]\).
\end{theorem}
\begin{proof}
First take the unique normalized representative. In Jacobi,
use the cyclic expression \(\sum_{\mathrm{cyc}}[[x,y],z]\)
as in Results I02--I03 of \cite{mg-lie}.
The coefficient in order \(t\) is zero since \(\phi\) is
a cocycle. At order \(t^2\), the contribution of the first-order
cochain on \((x,y,z)\) is
\[
 \begin{aligned}
   \sum_{\mathrm{cyc}}\phi(\phi(x,y),z)
    &= -e\,\phi(x,z)+c\,\phi(y,z)\\
    &= (cd-ea)x-(eb+ca)y .
 \end{aligned}
\]
If a \(t^2\)-cochain is added to the bracket, its contribution
at that order is the negative of its differential, by
Result I03. By Proposition~\ref{res:LD01}, every such differential
takes values in \(kz\) on \((x,y,z)\). It can cancel the
displayed vector only if both its coordinates are zero.

When both vanish, Jacobi for the bracket with only the linear
term is identically zero: its coefficients have degrees \(0,1,2\),
and each has vanished on an exterior basis for triples.
Changing the representative of the first-order class by a
coboundary preserves existence of either extension. The first-order
change is \(1+t f\), as in Result I02, and the same linear
operator is invertible modulo \(t^3\) and over \(k[[t]]\) by
inverting it as a formal power series. Transport the extension
by that change of basis.
\end{proof}

\begin{example}[Two extendable classes with an obstructed sum]
\label{res:LD03}
Let \(\phi_1(x,z)=x,\phi_1(y,z)=-y\), and
\(\phi_2(x,z)=0,\phi_2(y,z)=z\), with both cochains zero
on \((x,y)\). Each extends, but their sum is obstructed
already in degree two.
The extension defined by \(\phi_1\) specializes at \(t=0\)
to \(\mathfrak h\) and is isomorphic to
\(\mathfrak{sl}_2(k((t)))\) over the Laurent series field.
\end{example}
\begin{proof}
For \(\phi_1\) only \(a=1\) is nonzero, and for \(\phi_2\)
only \(e=1\) is nonzero. The equations of
Theorem~\ref{res:LD02} both hold in either case.
For their sum the obstruction vector is \(-x\), nonzero
in the quotient in Proposition~\ref{res:LD01}.
For the generic \(\phi_1\)-bracket put
\(H=-2z/t\), \(E=x\), \(F=-2y/t\). These form a basis over
\(k((t))\), and its brackets are \([H,E]=2E\),
\([H,F]=-2F\), and \([E,F]=H\), precisely the displayed
matrix-basis brackets of \(\mathfrak{sl}_2\).
\end{proof}

The general obstruction formalism for deformations over Artinian
bases is developed by Fialowski--Fuchs~\cite[Sections 1 and 3]{fialowski-fuchs}.
The calculation above concerns extensions of each first-order class
over one parameter.

\section{Coboundary Lie bialgebras}

\begin{definition}[Lie bialgebras and coboundary cobrackets]
\label{res:LD04}
For a finite-dimensional Lie algebra \(\mathfrak g\), put
\(x\cdot(y\wedge z)=[x,y]\wedge z+y\wedge[x,z]\).
A Lie bialgebra cobracket is a linear map
\(\delta:\mathfrak g\to\bigwedge^2\mathfrak g\)
such that
\[
 \delta([x,y])=x\cdot\delta(y)-y\cdot\delta(x)
\]
and the dual operation, given by
\(\langle[\alpha,\beta]_\delta,x\rangle
=\langle\alpha\wedge\beta,\delta(x)\rangle\), is a Lie bracket.
For \(r\in\bigwedge^2\mathfrak g\) the coboundary map
\(\delta_r(x)=x\cdot r\) always satisfies the first condition.
\end{definition}

\begin{proposition}[The invariant Schouten-square criterion]
\label{res:LD05}
Define a symmetric bracket on bivectors by bilinear extension of
\[
 \begin{aligned}
 [a\wedge b,c\wedge d]_{\mathrm S}
   ={}&[a,c]\wedge b\wedge d-[a,d]\wedge b\wedge c\\
      &-[b,c]\wedge a\wedge d+[b,d]\wedge a\wedge c .
 \end{aligned}
\]
The map \(\delta_r\) is a Lie bialgebra cobracket if and only if
\([r,r]_{\mathrm S}\) is fixed by the adjoint action
on \(\bigwedge^3\mathfrak g\).
\end{proposition}
\begin{proof}
The displayed formula is alternating in the two factors of
each bivector, so it descends to the exterior squares.
Interchanging the two bivectors changes the sign of the Lie
bracket and of the necessary permutation of the two remaining
vectors, making the formula symmetric. Applying
\(\operatorname{ad}_x\) to each term gives
\[
 x\cdot[r,s]_{\mathrm S}=[x\cdot r,s]_{\mathrm S}
                         +[r,x\cdot s]_{\mathrm S},
\]
using the Lie Jacobi identity in the bracketed factor.

For any \(\delta\), extend it in degree two by
\(\partial_\delta(y\wedge z)=\delta(y)\wedge z-y\wedge\delta(z)\).
Pairing \(\partial_\delta\delta(x)\) with \(\alpha\wedge\beta\wedge\gamma\)
gives the pairing of \(x\) with the cyclic sum
\(\sum_{\mathrm{cyc}}[[\alpha,\beta]_\delta,\gamma]_\delta\).
Indeed for a term \(u\wedge v\) of \(\delta(x)\) the pairing of
\(\delta(u)\wedge v-u\wedge\delta(v)\) is the cyclic sum of
\(
\langle[\alpha,\beta]_\delta,u\rangle\langle\gamma,v\rangle
-\langle[\alpha,\beta]_\delta,v\rangle\langle\gamma,u\rangle
\).
These pairings detect every trivector, so the dual Jacobi
identity is equivalent to \(\partial_\delta\delta=0\).

The bivector formula gives
\([r,y\wedge z]_{\mathrm S}=-\delta_r(y)\wedge z+
y\wedge\delta_r(z)\), first for a decomposable \(r\)
and then by linearity. Consequently
\[
 \partial_{\delta_r}\delta_r(x)
      =-[r,x\cdot r]_{\mathrm S}
      =-\tfrac12 x\cdot[r,r]_{\mathrm S}.
\]
The \(1\)-cocycle condition also holds because the adjoint
action is a representation. This proves the criterion.
\end{proof}

\begin{example}[Two cobrackets on the split rank-one algebra]
\label{res:LD06}
For the basis \(H,E,F\) with \([H,E]=2E\),
\([H,F]=-2F\), \([E,F]=H\), there are Lie bialgebra
structures from \(r_1=F\wedge E\) and \(r_2=H\wedge E\).
The first has
\(\delta_{r_1}(H)=0\), \(\delta_{r_1}(E)=H\wedge E\),
and \(\delta_{r_1}(F)=H\wedge F\).
The second has Schouten square zero and has
\(\delta_{r_2}(H)=2H\wedge E\), \(\delta_{r_2}(E)=0\),
and \(\delta_{r_2}(F)=-2E\wedge F\).
\end{example}
\begin{proof}
For \(r=u\wedge v\), the formula of
Proposition~\ref{res:LD05} gives
\([r,r]_{\mathrm S}=2[u,v]\wedge u\wedge v\).
For \(r_1\), this is a scalar multiple of \(H\wedge E\wedge F\).
The adjoint action of \(H\) on that trivector cancels the
weights \(2,-2\). The adjoint actions of \(E\) and \(F\)
each give repeated factors, so the trivector is invariant.
For \(r_2\), \([H,E]=2E\) gives a repeated factor in the square.
The criterion gives both bialgebras. Acting on the two
bivectors with the basis vectors gives the claimed formulas.
The first convention agrees with the simple-root normalization
of the preceding quantum chapter.
\end{proof}

\section{First-order changes of a bialgebra}

\begin{proposition}[Linearized associativity, coassociativity, and compatibility]
\label{res:LD07}
Let \(B\) be a bialgebra with multiplication \(m\), coproduct \(\Delta\),
unit \(1\), and counit \(\epsilon_B\).
Over \(k[s]/(s^2)\), write \(m_s=m+sf\), \(\Delta_s=\Delta+sg\)
and keep the unit and counit fixed. Let
\(J=1\otimes\tau\otimes1\) on \(B^{\otimes4}\), where
\(\tau\) interchanges the middle factors.
The bialgebra identities hold precisely when
\[
 \begin{aligned}
 &f(1,b)=f(b,1)=0,\quad \epsilon_B f=0,\quad
   g(1)=0,\quad (\epsilon_B\otimes1)g=(1\otimes\epsilon_B)g=0,\\
 &a f(b,c)-f(ab,c)+f(a,bc)-f(a,b)c=0,\\
 &(\Delta\otimes1)g+(g\otimes1)\Delta
        -(1\otimes\Delta)g-(1\otimes g)\Delta=0,\\
 &\Delta f+gm=(f\otimes m+m\otimes f)J(\Delta\otimes\Delta)
             +(m\otimes m)J(g\otimes\Delta+\Delta\otimes g).
 \end{aligned}
\]
Under pullback by \(1+s u\) with \(u(1)=0\) and
\(\epsilon_B u=0\), the two cochains change by
\[
 \begin{aligned}
 f&\longmapsto f+m(u\otimes1+1\otimes u)-u m,\\
 g&\longmapsto g+\Delta u-(u\otimes1+1\otimes u)\Delta.
 \end{aligned}
\]
\end{proposition}
\begin{proof}
Keeping the unit and counit supplies the first line, including
the condition that the coproduct is unital and the counit is
multiplicative. For associativity compare
\(m_s(m_s\otimes1)\) and \(m_s(1\otimes m_s)\);
their difference in degree one gives the second line up to
an overall sign. The degree-one part of
\((\Delta_s\otimes1)\Delta_s=(1\otimes\Delta_s)\Delta_s\)
is the third line. For compatibility, start from
\(\Delta_s m_s=(m_s\otimes m_s)J(\Delta_s\otimes\Delta_s)\).
On the left, the \(s\)-terms are \(\Delta f\) and \(gm\).
On the right one changes one of the four indicated operations
at a time. This gives the last line, and terms with two changes
are zero. These expansions exhaust the bialgebra identities.
In the pullback formulas use the inverse \(1-su\). Apply it
after multiplication and in both factors after comultiplication;
their linear terms are exactly the two displayed changes.
\end{proof}

These equations give the two operation components and change of
coordinates of the normalized bialgebra deformation complex.
For the general cohomology and higher brackets see
Markl~\cite[Section 6]{markl}.

\begin{theorem}[A nontrivial coproduct deformation of a polynomial Hopf algebra]
\label{res:LD08}
Let \(B=k[x,y]\) with \(x,y\) primitive. With multiplication
unchanged, the \(t\)-adically completed algebra \(B[[t]]\)
admits a Hopf structure extending the original one:
\[
 \Delta_t(x)=x\otimes1+1\otimes x,\qquad
 \Delta_t(y)=y\otimes1+\exp(tx)\otimes y .
\]
The counit still vanishes on \(x,y\). Its first-order class
as a bialgebra deformation is nonzero. On primitives of
degree one the linearized skew coproduct gives
\(\delta(x)=0\) and \(\delta(y)=x\wedge y\).
\end{theorem}
\begin{proof}
The formulas determine a continuous algebra map to
\((B\otimes B)[[t]]\), since the target is commutative and
both images have prescribed constant terms. As \(x\) is
primitive, the formal series \(a=\exp(tx)\) is group-like:
the exponential of the sum of the two commuting copies of
\(tx\) is the product of their exponentials.
Both iterated coproducts of \(y\) are
\[
 y\otimes1\otimes1+a\otimes y\otimes1+a\otimes a\otimes y,
\]
and they agree on \(x\) by primitivity. The counit
identities follow on generators. Define the continuous algebra
map \(S\) by \(S(x)=-x\), \(S(y)=-a^{-1}y\), so \(S(a)=a^{-1}\).
The two antipode identities hold on \(x\) and \(y\):
on \(y\) the sums are \(S(y)+a^{-1}y=0\) and
\(y+aS(y)=0\). The two sides are algebra maps by commutativity,
so the identities hold throughout the completion.

At first order the coefficient of \(t\) sends \(y\) to
\(x\otimes y\). The difference of this tensor and its flip
is \(x\otimes y-y\otimes x\), identified with \(x\wedge y\).
Since
\(\Delta\) is cocommutative, the second change of cochains
in Proposition~\ref{res:LD07} is symmetric under the flip,
for every \(u\). Thus the skew part cannot be removed by
a first-order change of coordinates. Its values on \(x,y\)
give the asserted cobracket on the abelian two-dimensional
Lie algebra they span. Its dual is a Lie bracket, since the
Jacobiator is alternating on a two-dimensional space.
\end{proof}

For the quantization problem with a prescribed Lie bialgebra
and its classical limit, see the definition and existence
theorem of Etingof--Kazhdan
~\cite[Section 3 and Theorem 9.6]{etingof-kazhdan}.

\backmatter
\begingroup
\footnotesize
\begin{thebibliography}{99}
\phantomsection
\addcontentsline{toc}{chapter}{Bibliography}
\renewcommand{\runningchapter}{Bibliography}
\setlength{\itemsep}{2pt}

\bibitem{etingof-kazhdan}
P. Etingof and D. Kazhdan.
\emph{Quantization of Lie Bialgebras, I}.
Selecta Mathematica, 1996. Archived as arXiv:q-alg/9506005v5.

\bibitem{etingof-semenyakin}
P. Etingof and M. Semenyakin.
\emph{A Brief Introduction to Quantum Groups}.
Lecture notes, archived author's version.

\bibitem{fialowski-fuchs}
A. Fialowski and D. Fuchs.
\emph{Construction of Miniversal Deformations of Lie Algebras}.
Journal of Functional Analysis, 1999. Archived as arXiv:math/0006117v1.

\bibitem{hernandez-affine}
D. Hernandez.
\emph{An Introduction to Affine Kac--Moody Algebras}.
CTQM Master Class Series 2, 2006.

\bibitem{jantzen-alggroup}
J.~C. Jantzen.
\emph{Representations of Algebraic Groups}.
Second edition, American Mathematical Society, 2007 reprint.

\bibitem{jantzen-modular}
J.~C. Jantzen.
\emph{Representations of Lie Algebras in Positive Characteristic}.
Archived author's notes.

\bibitem{kac}
V.~G. Kac.
\emph{Infinite-Dimensional Lie Algebras}.
Third edition, Cambridge University Press, 1990, corrected reprint, 1995.

\bibitem{kassel}
C. Kassel.
\emph{Quantum Groups}.
Graduate Texts in Mathematics 155, Springer, 1995.

\bibitem{markl}
M. Markl.
\emph{Intrinsic Brackets and the \(L_\infty\) Deformation Theory of Bialgebras}.
Archived as arXiv:math/0411456v7.

\bibitem{musson}
I.~M. Musson.
\emph{Lie Superalgebras and Enveloping Algebras}.
Graduate Studies in Mathematics 131, American Mathematical Society, 2012.

\bibitem{serganova}
V. Serganova.
\emph{Kac--Moody Superalgebras and Integrability}.
Archived author's manuscript.

\bibitem{stradefarnsteiner}
H. Strade and R. Farnsteiner.
\emph{Modular Lie Algebras and Their Representations}.
Monographs and Textbooks in Pure and Applied Mathematics 116, Dekker, 1988.

\bibitem{mg-algebra-core}
MathGenome.
\emph{Algebra Core}. P1-03, revision in MathGenome v0.25.0.

\bibitem{mg-lie}
MathGenome.
\emph{Lie Algebras and Finite-Dimensional Representations}.
ALG-08, 74-page revision, 2026.

\bibitem{mg-reptopics}
MathGenome.
\emph{Topics in Representation Theory}.
ALG-07, 37-page revision, 2026. Results LO09--LO14 concern untwisted affine algebras.

\bibitem{mg-rm03d}
MathGenome.
\emph{Topics in Rings, Modules, and Associative Algebras}.
ALG-03, 38-page revision, 2026. Chapter 3 includes \(U_q(\mathfrak{sl}_2)\).

\end{thebibliography}
\endgroup

\end{document}
